% Typesetting of Heinrich Weber's "Lehrbuch der Algebra, Band 3"
% Pages corresponding to PDF pages 251-300
% Content: Sections 67-76 on elliptic modular functions and transformation equations

\documentclass[12pt,a4paper]{article}

% Encoding and language
\usepackage[T1]{fontenc}
\usepackage[utf8]{inputenc}
\usepackage[ngerman]{babel}

% Mathematics packages
\usepackage{amsmath,amssymb,amsthm}
\usepackage{mathrsfs}
\usepackage{mathtools}

% Typography
\usepackage{microtype}
\usepackage{enumitem}
\usepackage{array}
\usepackage{booktabs}

% Page layout
\usepackage{geometry}
\geometry{a4paper,left=2.5cm,right=2.5cm,top=2.5cm,bottom=2.5cm}

% Headers and footers
\usepackage{fancyhdr}
\pagestyle{fancy}
\fancyhf{}
\fancyhead[C]{\small Siebenter Abschnitt}
\fancyfoot[C]{\thepage}
\renewcommand{\headrulewidth}{0.4pt}

% Section formatting
\renewcommand{\thesection}{\S\ \arabic{section}.}
\renewcommand{\thesubsection}{}
\renewcommand{\thesubsubsection}{}

% Custom commands for common notation
\newcommand{\fomega}{f(\omega)}
\newcommand{\fomeganull}{f(\omega_0)}
\newcommand{\fnull}{f_0}
\newcommand{\feins}{f_1}
\newcommand{\fzwei}{f_2}
\newcommand{\cnull}{c_0}
\newcommand{\ceins}{c_1}
\newcommand{\czwei}{c_2}
\newcommand{\cn}{c_n}
\newcommand{\const}{\text{const.}}

% Theorem-like environments
\theoremstyle{definition}
\newtheorem*{satz}{Satz}

% Title information
\title{Lehrbuch der Algebra --- Band 3, Teil 6}
\author{Heinrich Weber}
\date{}

\begin{document}

\maketitle
\thispagestyle{empty}

\vspace{1cm}
\begin{center}
\Large\textbf{Siebenter Abschnitt.}\\[0.5em]
\large\textbf{Elliptische Modulfunktionen und Transformationsgleichungen.}
\end{center}

\vspace{1cm}

\setcounter{section}{66}



\section{Besondere Transformationsgleichungen.}\label{sec:67}

Außerdem setzen wir noch:
\begin{equation}\label{eq:67-4}
P_{\nu,\mu}^{(n)} = e^{\frac{\nu(n^2-1)}{24n}\pi i} \prod_{h=0}^{n-1} \frac{\sigma\left(\frac{\nu\omega + \mu + 2h\tilde{\omega}}{n}\right)}{\sigma\left(\frac{\nu\omega + \mu}{n}\right) \cdot \sigma\left(\frac{2h\tilde{\omega}}{n}\right)} \quad [\S\ 67, (4)]
\end{equation}

Diese Größen $P$ lassen sich durch die Wurzeln der Teilungsgleichung ausdrücken, wenn man unter Anwendung von (3) die Quotienten
\[
\frac{P_{1,0}}{P_{0,0}}, \quad \frac{P_{0,1}}{P_{0,0}}, \quad \frac{P_{1,1}}{P_{0,0}}, \quad \frac{P_{\nu,\mu}}{P_{0,0}}
\]
bildet und vermittelst der Formeln des \S\ 42 elliptische Funktionen einführt.

Man erhält so
\begin{equation}\label{eq:67-5}
\left(\frac{P_{1,0}}{P_{0,0}}\right)^2 = -\frac{1}{2}\sqrt[n]{\frac{\Delta}{\Delta^{(n)}}} \cdot \frac{1}{\sqrt{1-\varkappa^2}} \cdot \frac{\sqrt{1-\varkappa^2\,\mathrm{sn}^2\frac{2K}{n}} + \mathrm{sn}\frac{2K}{n}\,\mathrm{dn}\frac{2K}{n}}{\sqrt{1-\varkappa^2\,\mathrm{sn}^2\frac{2K}{n}} - \mathrm{sn}\frac{2K}{n}\,\mathrm{dn}\frac{2K}{n}},
\end{equation}

\begin{equation}\label{eq:67-6}
\frac{P_{0,1}}{P_{0,0}} = \sqrt[n]{\frac{\Delta}{\Delta^{(n)}}} \cdot \frac{1}{\varkappa^2} \cdot \frac{\mathrm{sn}^2\frac{2K}{n}\,\mathrm{dn}^2\frac{2K}{n}}{1-\varkappa^2\,\mathrm{sn}^2\frac{2K}{n}}.
\end{equation}

Diese Ausdrücke zeigen nun, daß $P_{0,1}^2$, $P_{1,0}^2$, $P_{1,1}^2$ die Wurzeln von Transformationsgleichungen (Modulargleichungen) sind. $P_{0,1}^2$ ist die Wurzel einer Multiplikatorgleichung, und wenn $n$ eine Quadratzahl ist, so ist auch $P_{1,1}^2$ die Wurzel einer Multiplikatorgleichung.

Man kann in mannigfaltiger Weise die Funktionen $P^2$ miteinander kombinieren, um neue einfachere zu erhalten. Wir führen folgende an:
\begin{equation}\label{eq:67-7}
\begin{aligned}
P^{(n)} &= P_{0,1}^2 = e^{\frac{n^2-1}{12n}\pi i} \prod_{h=1}^{n-1} \frac{\sigma^2\frac{2h\tilde{\omega}}{n}}{\sigma^2\frac{h\tilde{\omega}}{n}}, \\
P^{(n)}_1 &= P_{1,0}^2 = e^{\frac{n^2-1}{12n}\pi i} \prod_{h=0}^{n-1} \frac{\sigma^2\frac{\omega+2h\tilde{\omega}}{n}}{\sigma^2\frac{h\tilde{\omega}}{n}}, \\
P^{(n)}_2 &= P_{1,1}^2 = e^{\frac{n^2-1}{12n}\pi i} \prod_{h=0}^{n-1} \frac{\sigma^2\frac{\omega+1+2h\tilde{\omega}}{n}}{\sigma^2\frac{h\tilde{\omega}}{n}}.
\end{aligned}
\end{equation}


\begin{equation}\label{eq:67-8}
\begin{aligned}
P^{(n)} &= e^{\frac{n^2-1}{12n}\pi i} \sqrt[n]{\frac{\Delta}{\Delta^{(n)}}} \cdot \frac{1}{\varkappa^{\frac{2(n-1)}{n}}} \cdot \prod_{h=1}^{n-1} \frac{\mathrm{sn}^2\frac{2hK}{n}\,\mathrm{dn}^2\frac{2hK}{n}}{1-\varkappa^2\,\mathrm{sn}^2\frac{2hK}{n}}, \\
P^{(n)}_1 &= e^{\frac{n^2-1}{12n}\pi i} \sqrt[n]{\frac{\Delta}{\Delta^{(n)}}} \cdot \frac{1}{\varkappa^{\frac{2(n-1)}{n}}} \cdot \prod_{h=0}^{n-1} \frac{\mathrm{sn}^2\frac{2hK}{n}\,\mathrm{cn}^2\frac{2hK}{n}}{1-\varkappa^2\,\mathrm{sn}^2\frac{2hK}{n}}, \\
P^{(n)}_2 &= e^{\frac{n^2-1}{12n}\pi i} \sqrt[n]{\frac{\Delta}{\Delta^{(n)}}} \cdot \frac{1}{\varkappa^{\frac{2(n-1)}{n}}} \cdot \prod_{h=0}^{n-1} \frac{\mathrm{sn}^2\frac{2hK}{n}\,\mathrm{dn}^2\frac{2hK}{n}}{1-\varkappa^2\,\mathrm{sn}^2\frac{2hK}{n}}.
\end{aligned}
\end{equation}

Ein wichtiges Resultat ergibt sich aber noch durch Anwendung des Multiplikationstheorems. Wenn man in der letzten Formel II des \S\ 57 $u = \frac{2h\tilde{\omega}}{n}$ setzt, so findet man
\begin{equation}\label{eq:67-9}
\prod_{h=0}^{n-1} \mathrm{sn}^2\frac{2h\tilde{\omega}}{n} = \frac{(-1)^{\frac{n-1}{2}}}{\varkappa^{n-1}} \cdot \frac{1}{P^{(n)}_1 \cdot P^{(n)}_2 \cdots P^{(n)}_{n-1}},
\end{equation}
worin, wie wir uns erinnern, $f_1$ eine ganze rationale Funktion ist, deren Koeffizienten rational aus $\varkappa^2$ und ganzen Zahlen gebildet sind.

Nehmen wir das Produkt aus diesen Ausdrücken, für $\nu = 1, 2, \ldots, n-1$, so folgt aus der letzten Gleichung (2) (weil bekanntlich $\sum h^2 = \frac{n(n-1)(2n-1)}{6}$):
\begin{equation}\label{eq:67-10}
P^{(n)} \cdot P^{(n)}_1 \cdot P^{(n)}_2 \cdots P^{(n)}_{n-1} = (-1)^{\frac{n-1}{2}} \cdot n^2 \cdot \left(\frac{\Delta^{(n)}}{\Delta}\right)^{\frac{n-1}{2n}}.
\end{equation}

Wir schließen hieraus, daß auch $P^{(n)}_1$ die Wurzel einer Transformationsgleichung ist. Dies ist nichts Neues, wenn $n$ durch 3 teilbar ist, wohl aber, wenn $n$ nicht durch 3, also $n^2-1$ durch 3 teilbar ist. Wir erhalten dann nämlich aus (10) mit Benutzung von (5), (7), (8):
\begin{equation}\label{eq:67-11}
P^{(n)}_1 = \left(\frac{P^{(n)}}{P^{(n)}_2}\right)^{\frac{n^2-1}{3}} \cdot \left(\frac{\mathrm{cn}^2\frac{2K}{n}\,\mathrm{dn}^2\frac{2K}{n}}{\mathrm{D}(\mathrm{sn}^2\frac{2K}{n})}\right)^{\frac{n^2-1}{3}},
\end{equation}

\begin{equation}\label{eq:67-12}
P^{(n)}_2 = \left(\frac{P^{(n)}}{P^{(n)}_1}\right)^{\frac{n^2-1}{3}} \cdot \left(\frac{\mathrm{sn}^2\frac{2K}{n}\,\mathrm{dn}^2\frac{2K}{n}}{\mathrm{D}(\mathrm{sn}^2\frac{2K}{n})}\right)^{\frac{n^2-1}{3}},
\end{equation}
wobei
\[
\mathrm{D}(\mathrm{sn}^2\tfrac{2K}{n}) = \prod_{h=1}^{n-1} (1 - \varkappa^2\,\mathrm{sn}^2\tfrac{2hK}{n}).
\]

Soll als Rationalitätsbereich der der rationalen Zahlen und rationalen Funktionen von $\varkappa^2$ aufrecht erhalten werden, so schreibt man die letzte Gleichung besser in der Form [vgl.\ \S\ 54, (4)]:
\begin{equation}\label{eq:67-13}
\left(P^{(n)}_1\right)^3 = \left(\frac{P^{(n)}}{P^{(n)}_2}\right)^{n^2-1} \cdot \frac{\mathrm{sn}^2\frac{2K}{n}\,\mathrm{cn}^2\frac{2K}{n}\,\mathrm{dn}^2\frac{2K}{n}}{\mathrm{D}(\mathrm{sn}^2\frac{2K}{n})^2} \cdot \frac{1}{(\mathrm{sn}^2\frac{2K}{n})^{\frac{2n^2+1}{3}} (\mathrm{cn}^2\frac{2K}{n})^{\frac{2n^2+1}{3}} (\mathrm{dn}^2\frac{2K}{n})^{\frac{n-1}{3}}}
\end{equation}
und schließt an auf den folgenden Satz:

\begin{satz}
Wenn $n$ nicht durch 3 teilbar ist, so sind $P^{(n)}$, $P^{(n)}_1$, $P^{(n)}_2$, $P^{(n)}_1 \cdot P^{(n)}_2 (P^{(n)})^{\frac{n-1}{3}}$, und wenn $n$ ein Quadrat ist, auch $P^{(n)}_1 \cdot P^{(n)}_2 (P^{(n)})^{\frac{n-1}{6}}$ Wurzeln von Transformationsgleichungen.
\end{satz}


\section{Zweite Darstellung der Wurzeln der Transformationsgleichungen.}\label{sec:68}

Wenn wir zunächst eine der Reihen ins Auge fassen, nämlich die, zu der die Wurzel $\varkappa$ gehört, also $u=1$, $\omega'=0$ setzen, so können wir auf (2), (3) des vorigen Paragraphen die Formeln (20), (21), (10), \S\ 34 anwenden. Setzt man dort, wenn $\nu$ ungerade ist, $n\alpha - \nu$ an Stelle von $\nu$ und benutzt die Formel
\[
\sqrt[n]{\frac{\sigma(n\alpha - \nu)}{\sigma(\nu)}} = (-1)^{\frac{g_1(g_1+1)}{2}} \cdot e^{\frac{\nu}{n}\eta_1(2n\alpha - \nu)} \cdot Dg_1g_2(\nu),
\]
dann kommen in diesen Formeln nur die geraden $\nu$ vor, und wenn man also $\nu = 2h$ setzt, so kann man die dortigen Formeln (20), (21) auch so schreiben:
\[
\begin{aligned}
\sum_{h=0}^{\frac{n-3}{2}} \frac{\psi_{2h}(\omega)}{\psi_0(\omega)} \cdot \frac{\sigma(2h\omega + \omega_0)}{\sigma(\omega_0)} &= \frac{1}{2} \cdot \frac{\psi_{n-1}(\omega)}{\psi_0(\omega)} \cdot \frac{\sigma(\omega_0)}{\sigma(n\omega + \omega_0)}, \\
\frac{\psi_{n-1}(\omega)}{\psi_0(\omega)} &= \sqrt[n]{\frac{\sigma(\omega_0)}{\sigma(n\omega + \omega_0)} \cdot H(\omega)}, \\
\frac{\psi_0(\omega)}{\psi_0(\omega)} &= \sqrt[n]{\frac{\sigma(\omega_0)}{\sigma(n\omega + \omega_0)}} \cdot \frac{1}{H(\omega)},
\end{aligned}
\]
worin $(-1)^{\frac{g_1}{2}} = (-1)^{\frac{g_1}{8}}$ und wenn wir also die dem Falle $u=1$, $\omega'=0$ entsprechenden Funktionen $P$ mit $P^{(n)}_0$ bezeichnen, so folgt aus \S\ 67, (2) und (4) mit Rücksicht auf \S\ 34, (10):
\begin{equation}\label{eq:68-1}
\begin{aligned}
P^{(n)}_0 &= \frac{\psi_0(\omega)}{\psi_0(\omega)} \cdot \frac{f(\omega_0)}{f(n\omega)}, \\
P^{(n)}_1 &= \frac{\psi_0(\omega)}{\psi_0(\omega)} \cdot \frac{f(\omega_0)}{f(n\omega)}, \\
P^{(n)}_2 &= \frac{\psi_0(\omega)}{\psi_0(\omega)} \cdot \frac{f(\omega_0)}{f(n\omega)}.
\end{aligned}
\end{equation}

Um durch solche Formeln die sämtlichen Wurzeln der Transformationsgleichungen darzustellen, bestimmen wir eine lineare Transformation folgendermaßen:
\begin{equation}\label{eq:68-3}
\omega = \frac{a\omega' + b}{c\omega' + d} \quad (a,b,c,d \text{ ganz}, \ ad-bc=1),
\end{equation}
\begin{equation}\label{eq:68-4}
\omega = u, \quad \omega' = v \pmod{n},
\end{equation}
und ersetzen $\omega$ in (1), (2) durch
\begin{equation}\label{eq:68-5}
\omega = \frac{\omega + v}{u} \quad \text{oder} \quad \frac{v\omega + 1}{u},
\end{equation}
auf die linke Seite lassen sich dann die Formeln (3) bis (6) \S\ 39 und (19), \S\ 38 anwenden und ergeben:
\begin{equation}\label{eq:68-6}
\frac{f_r(\frac{\omega + v}{u})}{f(\omega)} = \sqrt[n]{\frac{\Delta^{(n)}}{\Delta}} \cdot \frac{\psi_0(\frac{a\omega + b}{c\omega + d})}{\psi_0(\omega)} \cdot \frac{f(\frac{a\omega + b}{c\omega + d})}{f(\omega)},
\end{equation}
worin $\varepsilon$ die in \S\ 33, (15), (18) genau definierte 24ste Einheitswurzel ist.

Nun lassen sich zwei Transformationen
\[
\omega = \frac{a\omega + b}{c\omega + d}, \quad \omega = \frac{a'\omega + b'}{c'\omega + d'}
\]
von denen die erste linear, die andere von der $n$ten Ordnung ist, so bestimmen, daß
\begin{equation}\label{eq:68-7}
\begin{pmatrix} a & b \\ c & d \end{pmatrix} \cdot \begin{pmatrix} 1 & 0 \\ 0 & n \end{pmatrix} = \begin{pmatrix} \alpha & \beta \\ 0 & \theta \end{pmatrix} \cdot \begin{pmatrix} a_1 & b_1 \\ c_1 & d_1 \end{pmatrix}
\end{equation}

Denn schreiben wir die Relation (7) so:
\[
\begin{pmatrix} a & b \\ c & d \end{pmatrix} \cdot \begin{pmatrix} 1 & 0 \\ 0 & n \end{pmatrix} = \begin{pmatrix} \alpha & \beta \\ 0 & \theta \end{pmatrix} \cdot \begin{pmatrix} a_1 & b_1 \\ c_1 & d_1 \end{pmatrix},
\]
so leitet man daraus einfach her:
\begin{equation}\label{eq:68-8}
\theta = a\delta - c\beta, \quad \alpha = n, \quad \nu\alpha = \theta a_1, \quad \nu\beta = \theta b_1, \quad \mu = a - c\beta.
\end{equation}

Hierdurch ist zunächst $\theta$ bestimmt als der größte gemeinschaftliche Teiler von $\delta$ und $n$ oder [wegen (4)] von $v$ und $n$. Denn wäre $\theta$ nicht der größte gemeinschaftliche Teiler von $n = \theta\alpha$ und $\beta = \theta\beta'$, so müßten $\alpha$ und $\beta'$ einen gemeinsamen Teiler haben und folglich auch $\beta$ und $\alpha = \alpha' - c\beta$, was nicht sein kann. Dadurch ist zugleich $\alpha = n/\theta$ bestimmt, und $c$ muß der Kongruenz
\begin{equation}\label{eq:68-9}
\theta c \equiv \beta \pmod{n}
\end{equation}
genügen, wodurch jedoch $c$ nicht absolut, sondern nur nach dem Modul $\alpha$ bestimmt ist. Man kann also $c$ z.~B.\ noch an die Bedingung knüpfen, daß es durch irgend eine Potenz von 2, oder wenn $\alpha$ nicht durch 3 teilbar ist, durch irgend eine Potenz von 3 teilbar sein soll, was wir später bisweilen tun werden.

Die drei Zahlen $\alpha$, $\theta$, $c$ können keinen gemeinsamen Teiler haben, denn ein solcher wäre [nach (8)] auch Teiler von $\alpha$ und $\beta$, was unmöglich ist. Bezeichnen wir also den größten gemeinsamen Teiler von $\alpha$ und $c$ mit $e$, so muß $c$ relativ prim zu $e$ sein.

Die Zahlen $\alpha$, $\theta$, $c$, letztere modulo $\alpha$, sind wegen (4) und (8) durch die beiden Zahlen $u$, $u'$ völlig bestimmt und ändern sich nicht, wenn $u$, $u'$ mit einem gemeinsamen, zu $n$ teilerfremden Faktor multipliziert, d.~h.\ durch ein anderes Paar derselben Reihe ersetzt werden. Zerlegt man $n$ irgendwie in zwei Faktoren $\alpha$, $\theta$, so kann man für $c$ noch
\[
\frac{\alpha}{e} = \varphi(e)
\]
nach dem Modul $\alpha$ verschiedene zu $e$ teilerfremde Werte annehmen. Jede dieser Annahmen führt zu einem Zahlenpaar $u$, $u'$ durch die Kongruenzen
\begin{equation}\label{eq:68-10}
u \equiv \alpha' \pmod{\alpha}, \quad u' \equiv \beta' \pmod{\alpha},
\end{equation}
worin $\beta'$ eine beliebige, zu $\alpha$ teilerfremde Zahl bedeutet und $\alpha'$ so bestimmt wird, daß $\alpha'$, $\beta'$ relativ prim werden, und es ist auch leicht [nach \S\ 65, (3)] einzusehen, daß, so lange $\alpha$, $\theta$, $c$ dieselben bleiben, die nach (10) bestimmten Zahlenpaare $u$, $u'$ derselben Reihe angehören.

Wir nennen die Zahlen $\alpha$, $\theta$, $c$ (wie in \S\ 27) die \emph{Transformationszahlen} und $n$ den \emph{Transformationsgrad}.

Das Zahlensystem $\alpha$, $\theta$, $c$ ist also vollständig charakteristisch für eine Reihe, und die Anzahl der Reihen ist gleich der Anzahl dieser Zahlensysteme, woraus für die Zahl $\psi(n)$, (\S\ 63) folgt:
\begin{equation}\label{eq:68-11}
\psi(n) = n \sum_{\theta|n} \frac{\varphi(\theta)}{\theta},
\end{equation}
wenn die Summe auf alle Divisoren $\theta$ von $n$ erstreckt wird.

Es ist von Interesse, diese Relation auch direkt zu beweisen, wobei die Beschränkung auf ein ungerades $n$ wegfallen kann. Betrachten wir $\psi(n)$ jetzt als Zeichen für die Summe (11), so ergibt sich, wenn $n'$, $n''$ relativ prim sind, zunächst
\[
\psi(n')\psi(n'') = \psi(n'n''),
\]
und es bleibt also nur übrig, die Summe $\psi(n)$ für den Fall zu bestimmen, daß $n = p^\varrho$ eine Primzahlpotenz ist. In diesem Falle ist nun $e$ gleich dem kleineren der beiden Divisoren $\alpha$, $\theta$, und wenn $\alpha = \theta$ ist, $e = \alpha$. Wenn wir also das erste und letzte Glied der Summe $\psi(n)$ absondern, so erhalten wir
\[
\psi(p^\varrho) = 1 + \sum_{\substack{\alpha = p^a,\ \theta = p^b \\ 1 < a \leq \varrho}} \frac{\varphi(\theta)}{\theta} + p^\varrho.
\]

Es ist aber
\[
\sum_{a=1}^{\varrho-1} \frac{\varphi(p^a)}{p^a} = \sum_{a=1}^{\varrho-1} \frac{p^a - p^{a-1}}{p^a} = \varrho - 1 - \frac{\varrho - 1}{p},
\]
also
\[
\psi(p^\varrho) = 1 + p^\varrho + p^\varrho \sum_{a=1}^{\varrho-1} \frac{\varphi(p^a)}{p^a} = 1 + p^\varrho + p^\varrho \left(\varrho - 1 - \frac{\varrho - 1}{p}\right),
\]
woraus sich für $\psi(n)$ der Ausdruck
\[
\psi(n) = n \prod_{p|n} \left(1 + \frac{1}{p}\right)
\]
ergibt, wie in \S\ 63.\footnote{Dies folgt aus der Dirichletschen Formel für $\psi(n)$.}

Nach (7) ist nun in den Formeln (6) für $n\omega'$ zu setzen:
\[
n\omega' = \frac{\alpha\omega + \beta}{\theta}, \quad \omega = \frac{\theta\omega' - \beta}{\alpha}
\]
und es lassen sich die linearen Transformationen der $f$-Funktionen [\S\ 40, (4), (8), (11)] anwenden. Wir nehmen dabei
\begin{equation}\label{eq:68-12}
c \equiv 0 \pmod{16},
\end{equation}
so daß nach (3) und (8):
\[
\frac{a\omega + b}{c\omega + d} \equiv \frac{a\omega}{c\omega} \pmod{16}.
\]

Bezeichnen wir mit $\varepsilon$ die dritte Einheitswurzel
\begin{equation}\label{eq:68-13}
\varepsilon = e^{\frac{2\pi i}{3}},
\end{equation}
so erhalten wir
\[
\frac{f_r(\frac{a\omega + b}{c\omega + d})}{f(\omega)} = \varepsilon^{\frac{c + d}{3}} \cdot \frac{f(\frac{\alpha\omega + \beta}{\theta})}{f(\omega)}.
\]

Der Quotient
\[
\frac{f(\frac{\alpha\omega + \beta}{\theta})}{f(\omega)}
\]
gibt nach \S\ 54, (3), wenn man
\[
u(\omega) = \sqrt[4]{\varkappa}
\]
setzt, die Größen
\[
\sqrt[\theta]{\frac{\varkappa}{\varkappa'}}, \quad \sqrt[\theta]{\frac{\varkappa'}{\varkappa''}}, \quad \sqrt[\theta]{\frac{\varkappa''}{\varkappa}}
\]
als Wurzeln einer Modulargleichung, und dies ist die Jacobische Modulargleichung.\footnote{Vgl.\ Dedekind: \"Uber die elliptischen Modulfunktionen. Journal f.\ Mathematik, Bd.~83.}


Ist $n$ nicht durch 3 teilbar, so nehmen wir
\begin{equation}\label{eq:68-17}
c \equiv 0 \pmod{3}
\end{equation}
an, wodurch $\varepsilon$ den Wert 1 erhält.

In gleicher Weise kann man die Transformation der $\eta$-Funktion [\S\ 38, (15), (18), (19)] auf die letzte der Gleichungen (6) anwenden und erhält:
\begin{equation}\label{eq:68-18}
\frac{\eta_r(\frac{a\omega + b}{c\omega + d})}{\eta(\omega)} = \sqrt[\theta]{\frac{\varkappa\varkappa'}{\varkappa''}} \cdot \varepsilon^{\frac{c + d}{3}}
\end{equation}
(wobei es schon genügen würde, wenn $c$ durch 8 teilbar ist). Ist $n$ durch 3 nicht teilbar und $c$ durch 3 teilbar, so ist auch hierin $\varepsilon = 1$ zu setzen.

Die Bestimmung des Vorzeichens in (18) hat für uns nur in dem Falle Interesse, wo $n$ ein Quadrat ist. Es ist aber [nach (8)]
\[
\left(\frac{\alpha}{\theta}\right) = \left(\frac{\alpha}{\theta}\right) \cdot \left(\frac{\theta}{\alpha}\right),
\]
letzteres nach dem Reziprozitätsgesetz der quadratischen Reste, wel $\alpha \equiv 1 \pmod{4}$. Wenn nun $n$ ein Quadrat ist, so sind auch $\alpha:\theta$, $\beta:\theta$ Quadrate und es ergibt sich:
\[
\sqrt{\frac{\alpha}{\theta}} = \sqrt{\alpha} \cdot \sqrt{\frac{1}{\theta}},
\]
also wird in diesem Falle
\[
\frac{c + d}{3} = \frac{c + d}{3} \pmod{8}.
\]

\begin{satz}
Die zur Charakterisierung einer Reihe aufgestellten Formeln (10) sind ein spezieller Fall eines allgemeineren Systems, das man erhält, indem man auf $\alpha'$, $\beta'$ in (10) eine lineare Transformation anwendet. Man erhält dann folgendes:
\end{satz}

\footnotetext[2]{Ist $M$ der Jacobische Multiplikator, so ist}

Sind $a$, $b$, $c$, $d$ vier der Bedingung:
\[
ad - bc = n
\]
genügende ganze Zahlen ohne gemeinsamen Teiler, so erhält man die einer Reihe entsprechenden Zahlenpaare $u$, $u'$, wenn man in
\begin{equation}\label{eq:68-19}
u \equiv a\alpha' + b\beta' \pmod{n}, \quad u' \equiv c\alpha' + d\beta' \pmod{n}
\end{equation}
$\alpha'$, $\beta'$ alle und nur solche Werte durchlaufen läßt, bei denen $u$, $u'$ ohne gemeinsamen Teiler mit $n$ sind.

Daß zwei den Kongruenzen (19) entsprechende Wertpaare $u$, $u'$ wirklich derselben Reihe angehören, ergibt sich unmittelbar aus \S\ 65 (3), und ebenso ist selbstverständlich, daß alle Zahlenpaare einer Reihe in dieser Form enthalten sind, da man $\alpha'$, $\beta'$ durch $h\alpha'$, $h\beta'$ ersetzen kann, wenn $h$ relativ prim zu $n$ ist. Daß man sämtliche Reihen auf diesem Wege bekommt, zeigen die Formeln (10).


\section{Die Invariantengleichung.}\label{sec:69}

Unter den Transformationsgleichungen verdient eine ein besonderes Interesse, nämlich die, deren Wurzeln die $\psi(n)$ Größen
\begin{equation}\label{eq:69-1}
j\left(\frac{c + \omega}{\theta}\right)
\end{equation}
oder nach der Bestimmung (19), \S\ 68 die damit identischen Größen
\begin{equation}\label{eq:69-2}
j\left(\frac{a\omega + b}{c\omega + d}\right)
\end{equation}
sind, wenn $j(\omega)$ die in \S\ 46 definierte Invariante ist.

Diese Gleichung heißt die \emph{Invariantengleichung}.

Die Funktion $j(\omega)$ läßt sich nach \S\ 54 rational durch $f(\omega)^2$ darstellen, nämlich
\begin{equation}\label{eq:69-3}
j(\omega) = \frac{[f(\omega)^{16} - f_1(\omega)^8 f_2(\omega)^8]^3}{f(\omega)^8 f_1(\omega)^8 f_2(\omega)^8} = \frac{[f(\omega)^{16} + 16]^3}{f(\omega)^8 \cdot 16},
\end{equation}
so daß also nach den Resultaten des vorigen Paragraphen die Größen (1) die Wurzeln einer Gleichung sind, deren Koeffizienten rationale Funktionen von $\varkappa^2$ sind.

Setzt man aber für $f(\omega)$ in (3) eine der früher gefundenen Entwickelungen, z.~B.\ \S\ 24, a:
\[
f(\omega) = \sqrt[4]{2} \cdot q^{-\frac{1}{24}} \prod_{m=1}^{\infty} (1 + q^{2m-1}),
\]
so erkennt man, daß $j(\omega)$ sich in eine nach Potenzen von $q$ fortschreitende Reihe entwickeln läßt, welche die Form hat
\begin{equation}\label{eq:69-4}
j(\omega) = q^{-2} + a_0 + a_1 q^2 + a_2 q^4 + \cdots,
\end{equation}
worin die $a_0$, $a_1$, $a_2$, \ldots\ ganze rationale Zahlen sind, die sich successive berechnen lassen (es ergibt sich z.~B.\ $a_0 = 744$, $a_1 = 196\,884$). Die Entwickelungen der Größen (1) beginnen also mit
\begin{equation}\label{eq:69-5}
q^{-\frac{2}{n}} e^{\frac{2\pi i h}{n}} \quad (h = 0, 1, \ldots, n-1)
\end{equation}
und sind daher alle voneinander verschieden. Die Invariantengleichung ist also nach \S\ 65 irreducibel.

Es gibt aber einen zweiten Weg, um zu dieser wie überhaupt zu den Transformationsgleichungen zu gelangen, den wir jetzt zunächst bei der Invariantengleichung kennen lernen wollen. Diese Ableitung stützt sich auf die Sätze des \S\ 54 über Modulfunktionen und gilt auch für ein gerades $n$. Hier ist es der Satz 4, \S\ 54, der zur Anwendung kommt:

\begin{satz}[I.]
Jede Modulfunktion, die durch die beiden Transformationen
\begin{equation}\label{eq:69-6}
\omega' = \omega + 1, \quad \omega' = -\frac{1}{\omega}
\end{equation}
\begin{equation}\label{eq:69-7}
\omega' = \frac{a\omega + b}{c\omega + d}, \quad ad - bc = n
\end{equation}
ungeändert bleibt, ist eine rationale Funktion von $j(\omega)$.
\end{satz}

Wir weisen zunächst nach, daß durch Anwendung der Substitutionen (6), (7) die $\nu$ Größen (1) untereinander vertauscht werden. Wir haben die Zusammensetzung
\begin{equation}\label{eq:69-8}
\begin{pmatrix} a & b \\ c & d \end{pmatrix} \cdot \begin{pmatrix} 1 & 0 \\ 0 & n \end{pmatrix} = \begin{pmatrix} a & n\beta \\ \gamma & n\delta \end{pmatrix} \cdot \begin{pmatrix} \alpha & \beta \\ \gamma & \delta \end{pmatrix},
\end{equation}
wenn
\begin{equation}\label{eq:69-9}
G = c + d \equiv 0 \pmod{n}
\end{equation}
gesetzt wird, und es geht durch die Transformation (6), da $j(\omega)$ durch jede lineare Transformation ungeändert bleibt,
\[
j\left(\frac{a\omega + b}{c\omega + d}\right) \text{ in } j\left(\frac{a\omega + b + n\beta}{c\omega + d + n\delta}\right)
\]
über.

Ferner bestimmen wir $\alpha_1$, $\theta_1$, $c_1$, so daß
\[
\begin{pmatrix} a_1 & b_1 \\ c_1 & d_1 \end{pmatrix} \cdot \begin{pmatrix} \alpha & \beta \\ 0 & \theta \end{pmatrix} = \begin{pmatrix} a & b \\ c & d \end{pmatrix} \cdot \begin{pmatrix} 1 & 0 \\ 0 & n \end{pmatrix}.
\]

Dazu ist erforderlich
\[
\alpha_1 a = a, \quad \alpha_1 \beta + \theta_1 b_1 = b, \quad c_1 a = c, \quad c_1 \beta + \theta_1 d_1 = d,
\]
also
\[
\alpha_1 = 1, \quad \theta_1 = n, \quad c_1 = c, \quad d_1 = d.
\]

Es ist also $\theta_1$ bestimmt als der größte gemeinschaftliche Teiler von $\alpha$ und $c$, und damit zugleich, wegen $\alpha_1 d_1 - b_1 c_1 = n$, auch $\alpha_1$.

Nach den beiden Gleichungen (12) kennt man jetzt $\beta_1$, $\delta_1$ als relative Primzahlen und kann $\alpha_1$, $\gamma_1$ aus der diophantischen Gleichung
\[
\alpha_1 \delta_1 - \beta_1 \gamma_1 = 1
\]
bestimmen. Ist dies geschehen, so folgt aus den beiden Gleichungen (11)
\[
G = \alpha_1 \beta_1, \quad n = -\alpha_1 \gamma_1,
\]
wodurch auch $c_1$ bestimmt ist, und es zeigt sich zugleich, daß $\theta_1$ der größte gemeinschaftliche Teiler von $a_1$, $c_1$ ist. Ersetzt man $\alpha_1$, $\gamma_1$ durch eine andere Lösung $\alpha_1 + h\beta_1$, $\gamma_1 + h\delta_1$, so ändert sich nur $c_1$ um ein Vielfaches von $a_1$.

Durch die Transformation (7) geht also
\[
j\left(\frac{c + \omega}{\theta}\right) \text{ in } j\left(\frac{c_1 + \omega}{\theta_1}\right) = j\left(\frac{a_1 \omega + b_1}{c_1 \omega + d_1}\right)
\]
über.

Bilden wir nun eine symmetrische Funktion der sämtlichen Größen (1), etwa für ein unbestimmtes $x$ das Produkt
\[
F_n[x, j(\omega)] = \prod_{h=1}^{\psi(n)} \left(x - j\left(\frac{a_h \omega + b_h}{c_h \omega + d_h}\right)\right),
\]
so ändert sich diese Funktion nicht durch die Transformationen (6), (7) und ist also nach dem oben erwähnten Satz eine rationale Funktion von $j(\omega)$. Außerdem ist sie eine ganze rationale Funktion $\nu$ten Grades von $x$ mit dem Anfangsglied $x^\nu$, und wir bezeichnen sie mit $F_n[x, j(\omega)]$. Die Gleichung
\begin{equation}\label{eq:69-14}
F_n[x, j(\omega)] = 0
\end{equation}
hat die Größen $j\left(\frac{c + \omega}{\theta}\right)$ zu Wurzeln und ist die \emph{Invariantengleichung}, deren wichtigste Eigenschaften wir nun ableiten wollen.

\textbf{1.} Die Invariantengleichung ist irreducibel, wenn als Rationalitätsbereich der Inbegriff aller rationalen Funktionen von $j(\omega)$ mit beliebigen Zahlenkoeffizienten betrachtet wird. Denn besteht irgend eine rationale Gleichung
\begin{equation}\label{eq:69-15}
\Phi[j(n\omega),\ j(\omega)] = 0,
\end{equation}
so darf darauf jede beliebige lineare Transformation
\[
\omega' = \frac{\alpha\omega + \beta}{\gamma\omega + \delta}
\]
angewandt werden, und nach \S\ 29, (5) lassen sich, wenn
\[
\omega' = \frac{a\omega + b}{c\omega + d}
\]
eine beliebige Transformation $n$ter Ordnung ist, die linearen Transformationen
\[
\begin{pmatrix} \alpha & \beta \\ \gamma & \delta \end{pmatrix}, \quad \begin{pmatrix} \alpha_1 & \beta_1 \\ \gamma_1 & \delta_1 \end{pmatrix}
\]
immer so bestimmen, daß
\[
\begin{pmatrix} \alpha & \beta \\ \gamma & \delta \end{pmatrix} \cdot \begin{pmatrix} a & b \\ c & d \end{pmatrix} = \begin{pmatrix} a_1 & b_1 \\ c_1 & d_1 \end{pmatrix} \cdot \begin{pmatrix} 1 & 0 \\ 0 & n \end{pmatrix} \cdot \begin{pmatrix} \alpha_1 & \beta_1 \\ \gamma_1 & \delta_1 \end{pmatrix}.
\]

Daraus folgt aber, daß die Gleichung (9) durch jede der Größen
\[
j\left(\frac{c + \omega}{\theta}\right), \quad j\left(\frac{c_1 + \omega}{\theta_1}\right), \quad \ldots
\]
und mithin auch durch jede der Größen (1) befriedigt ist, woraus [wegen der Verschiedenheit der Größen (1)] die Irreducibilität von (14) folgt.

\textbf{2.} Die Funktion
\begin{equation}\label{eq:69-17}
F_n[x, j(\omega)] = \prod_{h=1}^{\psi(n)} \left(x - j\left(\frac{a_h \omega + b_h}{c_h \omega + d_h}\right)\right)
\end{equation}
hat für jeden endlichen Wert von $\omega$ mit positiv imaginärem Bestandteil, also auch für jeden endlichen Wert von $j(\omega)$, einen endlichen Wert und ist sonach eine ganze rationale Funktion von $j(\omega)$.

Suchen wir ferner nach (3) das Anfangsglied der Entwickelung der Funktion (17) nach Potenzen von $q$, so ergibt sich
\[
\prod_{h=1}^{\psi(n)} \left(x - q^{-\frac{2}{n}} e^{\frac{2\pi i h}{n}}\right) = x^{\psi(n)} + \cdots + (-1)^{\psi(n)} q^{-2\psi(n)/n} \cdot C,
\]
wenn $C$ eine endliche Konstante ist. Es ist daher
\[
\frac{F_n[x, j(\omega)]}{j(\omega)^{\psi(n)/n}}
\]
für ein unendliches $j(\omega)$ endlich, d.~h.\ der Grad von $F_n[x, j(\omega)]$ in bezug auf $j(\omega)$ ist ebenfalls der $\nu$te.

\textbf{3.} Es ist
\[
F_n[j(n\omega),\ j(\omega)] = 0
\]
und wenn wir $n\omega$ durch $\omega$ ersetzen:
\[
F_n[j(\omega),\ j(\tfrac{\omega}{n})] = 0.
\]

Da nun $j(\frac{\omega}{n})$ gleichfalls eine Wurzel der Invariantengleichung ist, so folgt, daß die beiden Gleichungen $\nu$ten Grades
\[
F_n[x, j(\omega)] = 0, \quad F_n[j(\omega), x] = 0
\]
eine Wurzel gemeinsam haben, und folglich, wegen der Irreducibilität der ersteren und der Gleichheit des Grades, alle. Es ist daher
\[
F_n(x, y) = C \cdot F_n(y, x)
\]
für unbestimmte Werte der Variablen $x$, $y$ und ein konstantes $C$. Daher auch, durch Vertauschung von $x$ mit $y$:
\[
F_n(y, x) = C \cdot F_n(x, y),
\]
also $C^2 = 1$ oder
\[
F_n(x, y) = C \cdot F_n(y, x).
\]

Wenn wir nun $x = y$ setzen, so folgt:
\[
F_n(y, y) = \pm F_n(y, y),
\]
also, wenn das untere Zeichen gilt,
\[
F_n(y, y) = -F_n(y, y) = 0.
\]

Dies ist aber nur möglich, wenn $n = 1$ ist [wo $F(x, y) = x - y$ wird], da sonst $F_n(x, y)$ durch $x - y$ teilbar sein müßte, was der Irreducibilität widerspricht. Daher ist immer, sobald $n > 1$ ist:
\begin{equation}\label{eq:69-18}
F_n(x, y) = F_n(y, x).
\end{equation}

\textbf{4.} Es sei
\[
\nu = \psi(n) = \psi(n')\psi(n''), \quad n = n'n'',
\]
und $n'$, $n''$ ohne gemeinsamen Teiler, und $\xi_1$, $\xi_2$, \ldots, $\xi_\nu$ seien die Wurzeln der Gleichung
\begin{equation}\label{eq:69-19}
F_{n'}[x, j(\omega)] = 0.
\end{equation}

Das Produkt
\begin{equation}\label{eq:69-20}
P = \prod_{h=1}^{\psi(n'')} F_{n''}[x, \xi_h]
\end{equation}
dessen Grad in bezug auf $x$ gleich
\[
\nu = \psi(n) = \psi(n')\psi(n'')
\]
ist, hängt als symmetrische Funktion der Wurzeln von (19), rational von $j(\omega)$ ab und verschwindet für
\[
\xi_h = j\left(\frac{c' + \omega}{\theta'}\right) = j\left(\frac{c + \omega}{\theta}\right) \quad [\S\ 69]
\]
d.~h.\ für eine Wurzel der Gleichung $F_n[x, j(\omega)] = 0$. Wegen der Irreducibilität der letzteren Gleichung, der Gleichheit des Grades und des Koeffizienten von $x^\nu$ ist also
\begin{equation}\label{eq:69-21}
F_n[x, j(\omega)] = P = \prod_{h=1}^{\psi(n'')} F_{n''}[x, \xi_h].
\end{equation}

\textbf{5.} Ist $n = p^\varrho$ eine Primzahlpotenz, so ist der Grad der Funktion $F_n[x, j(\omega)]$ gleich $p^{\varrho-1}(p+1)$, und die Gleichung
\begin{equation}\label{eq:69-22}
F_{p^{\varrho-1}}[x, j(\omega)] = 0
\end{equation}
ist vom Grade
\[
\psi(p^{\varrho-1}) = p^{\varrho-2}(p+1).
\]
Wir bezeichnen ihre Wurzeln mit $x_1$, $x_2$, \ldots, $x_r$.

Das Produkt
\[
P = \prod_{h=1}^{r} F_p[x, x_h]
\]
ist in bezug auf $x$ vom Grade $\nu(p+1) = p^{\varrho-2}(p+1)^2$; es hängt symmetrisch von den Wurzeln von (22), also rational von $j(\omega)$ ab und verschwindet für
\[
x_h = j\left(\frac{c + \omega}{\theta}\right) = j\left(\frac{c + \omega}{p\theta}\right) \quad [\S\ 69].
\]

Daher ist $P$ durch $F_n[x, j(\omega)]$ teilbar.

Aber der Grad von $P$ ist höher als der von $F_n$. Nehmen wir
\[
x_h = j\left(\frac{c + \omega}{p^{\varrho-1}\theta}\right), \quad c = 0, 1, 2, \ldots, p-1,
\]
wo $c$ jeden der Werte $0, 1, 2, \ldots, p-1$ haben kann, so verschwindet $F_p[x, x_h]$ für
\[
x = j\left(\frac{c' + \omega}{p\theta}\right) \quad [\S\ 69],
\]
d.~h.\ es haben $p$ von den Faktoren von $P$ einen bestimmten Faktor mit $F_{p^{\varrho-2}}[x, j(\omega)]$ gemein; daraus folgt, daß $P$ durch
\[
\{F_{p^{\varrho-2}}[x, j(\omega)]\}^p
\]
teilbar ist, und mithin, wie die Vergleichung der Grade und höchsten Glieder lehrt:
\begin{equation}\label{eq:69-23}
\prod_{h=1}^{r} F_p[x, x_h] = F_n[x, j(\omega)] \cdot \{F_{p^{\varrho-2}}[x, j(\omega)]\}^{p-1}.
\end{equation}

Diese Formel ist einer Ausnahme unterworfen für $p = 2$, weil in diesem Falle der Grad der Funktion auf der rechten Seite noch zu hoch ist. Für diesen Fall hat aber jeder der $p-1$ Faktoren des Produktes $P$ den Teiler $x - j(\omega)$, weil eben jedes $x_h$ Wurzel der Gleichung $F_p[x, j(\omega)]$ ist, und es tritt an Stelle von (23) die Formel
\begin{equation}\label{eq:69-24}
\prod_{h=1}^{r} F_2[x, x_h] = F_n[x, j(\omega)] \cdot \{F_{p^{\varrho-2}}[x, j(\omega)]\}^{p-1} \cdot (x - j(\omega))^{p-1}.
\end{equation}

Hierdurch ist die Lösung der Invariantengleichung $F_n = 0$ auf die successive Lösung solcher Fälle zurückgeführt, in denen $n$ eine Primzahl ist.

\textbf{6.} Während bei der Ableitung der Transformationsgleichungen aus den Teilungsgleichungen von Haus aus feststeht, daß die numerischen Koeffizienten in diesen Gleichungen rationale Zahlen sind, so lehrt uns die zweite Ableitung zunächst nichts über die Zahlenkoeffizienten. Wir können aber nachträglich beweisen, daß diese Koeffizienten nicht nur rationale, sondern auch ganze Zahlen sind und gelangen zugleich zu einem wichtigen Satz über die Teilbarkeit dieser Koeffizienten.

Wenn wir beweisen können, daß, wenn $p$ eine Primzahl ist, die Koeffizienten in $F_p(x, y)$ ganze Zahlen sind, so folgt das Gleiche aus 4.\ und 5.\ für jedes zusammengesetzte $n$.

Nach (4) ist $j(\omega)$ in eine Reihe von der Form entwickelbar
\begin{equation}\label{eq:69-25}
j(\omega) = q^{-2} \sum_{h=0}^{\infty} a_h q^{2h},
\end{equation}
deren Koeffizienten $a_h$ ganze Zahlen sind, und zwar $a_0 = 1$.

Bilden wir hiervon, wenn $p$ eine Primzahl ist, die $p$te Potenz, und beachten den für jede ganze Zahl gültigen Fermatschen Satz:
\[
a^p \equiv a \pmod{p},
\]
so folgt
\begin{equation}\label{eq:69-26}
j(\omega)^p = q^{-2p} \sum_{h=0}^{\infty} a_h q^{2hp} + p \sum_{h=1}^{\infty} d_h q^{2h},
\end{equation}
worin $d_h$ ebenfalls ganze Zahlen sind. Andererseits ist, wenn man in (25) $\omega$ durch $p\omega$ ersetzt:
\begin{equation}\label{eq:69-27}
j(p\omega) = q^{-2p} \sum_{h=0}^{\infty} a_h q^{2hp},
\end{equation}
und daraus, wenn man
\[
u = q^{-2}
\]
setzt:
\[
u - u = q^{-2} - q^{-2p} \sum_{h=0}^{\infty} a_h q^{2hp} = p \sum_{h=1}^{\infty} D_h q^{2h},
\]
wofür wir auch schreiben können:
\begin{equation}\label{eq:69-28}
(\nu - u)(\nu - u^p) = p \cdot q^{-2(p+1)} \sum_{h=0}^{\infty} c_h q^{2h},
\end{equation}
wenn $c_h$ ein drittes System ganzer Zahlen bedeutet.

Nun kommen in $F_p(x, y)$ die in bezug auf $x$, $y$ höchsten Glieder $x^{p+1}$, $y^{p+1}$ vor, und wir setzen demnach
\begin{equation}\label{eq:69-29}
F_p(x, y) = \sum_{h,k=0}^{p+1} c_{h,k} x^h y^k = \sum_{h,k=0}^{p+1} C_{h,k} \nu^h u^k,
\end{equation}
worin $c_{h,k}$ die zu bestimmenden Koeffizienten sind, die nach 3.\ der Bedingung
\[
c_{h,k} = c_{p+1-h,\,p+1-k} = C_{h,k}
\]
genügen. Um sie zu bestimmen, setzen wir in (29)
\[
\nu = u = j(\omega),
\]
wodurch $F_p$ verschwindet, und erhalten
\begin{equation}\label{eq:69-30}
(\nu - u)(\nu - u^p) = \sum_{h,k=0}^{p+1} C_{h,k} \nu^h u^k.
\end{equation}

Hieraus folgt zunächst, daß
\[
C_{p+1,0} = C_{0,\,p+1} = 1
\]
sein muß, da nach (28) bei der Entwickelung (der linken Seite) nach Potenzen von $q$ die Potenz $q^{-2(p+1)}$ nicht vorkommen kann, und wir können (29) jetzt in die Form setzen:
\begin{equation}\label{eq:69-31}
\begin{aligned}
(\nu - u)(\nu - u^p) &= \sum_{h,k=0}^{p} C_{h,k} (\nu^h u^k + \nu^k u^h) \\
&\quad - \sum_{h=0}^{p+1} C_{h,\,p+1-h} \nu^h u^{p+1-h}.
\end{aligned}
\end{equation}

Hierin sind die aus (25) und (27) sich ergebenden Entwickelungen von
\[
\nu^h u^k - \nu^k u^h, \quad \nu^h u^k + \nu^k u^h
\]
nach Potenzen von $q$ einzusetzen, deren Anfangsglieder
\[
q^{-2(h+k)}, \quad q^{-2(h+k)}
\]
die Koeffizienten 1 haben.

Auf der rechten Seite von (31) kommen nicht zwei Glieder vor, deren Entwickelung mit derselben Potenz von $q$ beginnt, denn aus
\[
hp + k = h'p + k'
\]
folgt $k \equiv k' \pmod{p}$ und mithin, da $k$ und $k' < p$ sind, $k = k'$, also $h = h'$.

Ordnet man daher die Reihen, welche die beiden Seiten von (31) darstellen, nach aufsteigenden Potenzen von $q$, und setzt dann die Koeffizienten gleich hoher Potenzen einander gleich, so erhält man eine Reihe linearer Gleichungen für die Unbekannten $C_{h,k}$, von denen jede folgende nur eine neue Unbekannte enthält, und diese mit dem Koeffizienten 1. Die aus der linken Seite sich ergebenden bekannten Glieder dieser Gleichungen sind nach (28) lauter durch $p$ teilbare ganze Zahlen, und es ergeben sich also für die $c_{h,k}$ ebenfalls ganzzahlige, durch $p$ teilbare Zahlwerte. Demnach haben wir
\begin{equation}\label{eq:69-32}
F_n(x,y) = (x - y)(x^p - y^p) - p \sum_{h,k=0}^{p+1} a_{h,k} x^h y^k,
\end{equation}
worin $a_{h,k}$ ganze Zahlen sind, die den Bedingungen
\[
a_{h,k} = a_{p+1-h,\,p+1-k}, \quad a_{p,0} = a_{0,\,p} = 0
\]
genügen.


\section{Transformationsgleichungen erster Stufe.}\label{sec:70}

Die Invariantengleichung ist von großer theoretischer Wichtigkeit teils wegen ihrer allgemeinen Gültigkeit (auch für gerade $n$), teils wegen der Leichtigkeit, mit der die linearen Transformationen auf $j(\omega)$ angewandt werden können. Die wirkliche Berechnung dieser Gleichung aber zeigt sich so kompliziert, und die Zahlenkoeffizienten sind so groß, daß die Berechnung bis jetzt nur in dem einfachsten Falle $p = 2$ durchgeführt ist. Dagegen kann man, indem man andere Modulfunktionen benutzt, weit einfachere Transformationsgleichungen erhalten.

Über das hierbei anzuwendende Prinzip bemerken wir folgendes:

Wenn irgend ein System von $\nu$ Funktionen von $\omega$ vorliegt, entsprechend den $\nu$ Systemen von Transformationszahlen $a$, $c$, $\theta$, die wir mit
\begin{equation}\label{eq:70-1}
\Phi_{a,c,\theta}(\omega)
\end{equation}
bezeichnen wollen und die isomorph mit den Funktionen
\begin{equation}\label{eq:70-2}
j\left(\frac{c + \omega}{\theta}\right)
\end{equation}
durch die Substitutionen
\begin{equation}\label{eq:70-3}
\omega' = \omega + 1, \quad \omega' = -\frac{1}{\omega}
\end{equation}
untereinander permutiert werden, so ist (1) rational durch (2) und durch $j(\omega)$ ausdrückbar, denn die Funktion
\begin{equation}\label{eq:70-4}
F_n[x, j(\omega)] = \prod_{h=1}^{\psi(n)} \left(x - \Phi_{a_h,c_h,\theta_h}(\omega)\right)
\end{equation}
bleibt durch die Substitutionen (3) ungeändert, und ist daher (für ein unbestimmtes $x$) eine rationale Funktion von $j(\omega)$ und überdies eine ganze rationale Funktion von $x$, höchstens vom Grade $\nu - 1$; $F_n[x, j(\omega)]$ hat dieselbe Bedeutung, wie in (17) des vorigen Paragraphen, und $\nu = \psi(n)$ ist der Grad dieser Funktion in bezug auf $x$. Aus (4) aber erhält man, indem man
\[
\Phi_{a,c,\theta}(\omega) = \frac{F_n[x, j(\omega)]}{\prod' (x - \Phi_{a_h,c_h,\theta_h}(\omega))}
\]
setzt:
\begin{equation}\label{eq:70-5}
\Phi_{a,c,\theta}(\omega) = \frac{F_n[x, j(\omega)]}{\Phi'[x, j(\omega)]}.
\end{equation}

Es gehört also $\Phi_{a,c,\theta}$ dem algebraischen Körper an, der aus den rationalen Funktionen von $j(\frac{c + \omega}{\theta})$ und $j(\omega)$ besteht, und wir können die Sätze von Bd.~I, \S\ 151 anwenden. Aus diesen folgt, daß die $\nu$ Größen $\Phi_{a_h,c_h,\theta_h}$ die Wurzeln einer Gleichung $\nu$ten Grades sind, deren Koeffizienten rational von $j(\omega)$ abhängen. Wenn die $\nu$ Größen $\Phi_{a_h,c_h,\theta_h}$ voneinander verschieden sind, so ist diese Gleichung irreducibel. Eine solche Gleichung nennen wir eine zum Transformationsgrad $n$ gehörige \emph{Transformationsgleichung erster Stufe}{}\footnote{In der ersten Auflage habe ich diese Gleichungen ``invariante Transformationsgleichungen'' genannt. Dieser Ausdruck ist von Klein beanstandet worden (Vorlesungen über ausgewählte Kapitel der Zahlentheorie, autographiertes Heft, Göttingen 1897). Ich schließe mich Kleins Vorschlag an, indem ich diese Gleichungen jetzt ``Transformationsgleichungen erster Stufe'' nenne, ohne hier näher auf die Begründung dieses Ausdruckes einzugehen.}. Jede andere Größe des Körpers kann durch ein solches $\Phi_{a,c,\theta}$ und durch $j(\omega)$ rational ausgedrückt werden.

Haben die Funktionen $\Phi_{a_h,c_h,\theta_h}$ die Eigenschaft, für jeden endlichen Wert von $\omega$ mit positivem, imaginärem Teil, also für jedes endliche $j(\omega)$ endlich zu bleiben, so ist die Funktion $F[x, j(\omega)]$ in (4) auch in bezug auf $j(\omega)$ eine ganze Funktion. Die Formel (5) könnte daher nur für solche besondere Werte von $\omega$ versagen, für die zwei Wurzeln der Invariantengleichung einander gleich werden.

Wenn die $\nu$ Größen $\Phi_{a_h,c_h,\theta_h}$ nicht alle voneinander verschieden sind, so zerfallen sie in Reihen von gleich vielen untereinander gleichen, und die aus (5) abzuleitende Gleichung $\nu$ten Grades ist eine Potenz einer irreducibeln Gleichung, die wir gelegentlich wohl auch als Transformationsgleichung bezeichnen werden (Bd.~I, \S\ 151, 2).

Sind die Größen $\Phi_{a_h,c_h,\theta_h}$ so gewählt, daß sie für kein endliches $\omega$ mit positiv imaginärem Bestandteil unendlich werden, so bleiben sie für jedes endliche $j(\omega)$ endlich, woraus folgt, daß die Koeffizienten in der Funktion des $\nu$ten Grades
\begin{equation}\label{eq:70-6}
\prod_{h=1}^{\psi(n)} (x - \Phi_{a_h,c_h,\theta_h})
\end{equation}
ganze rationale Funktionen von $j(\omega)$ sind, d.~h.\ die $\Phi_{a_h,c_h,\theta_h}$ sind ganze algebraische Funktionen von $j(\omega)$.

Richtet man die Funktionen $\Phi_{a_h,c_h,\theta_h}$ etwa durch geeignete Bestimmung von Konstanten, die darin noch verfügbar sind, so ein, daß sie auch für ein unendliches imaginäres $\omega$, d.~h.\ für $q = 0$ endlich bleiben, so sind diese Funktionen auch für ein unendliches $j(\omega)$ endlich, und die Koeffizienten in (6) sind konstant. Dies ist aber nur dadurch möglich, daß die $\Phi_{a_h,c_h,\theta_h}$ alle einer und derselben Konstanten $C$ gleich sind. Kennt man $\Phi_{a,c,\theta}$ als rationale Funktion einer anderen Größe $\Psi_{a,c,\theta}$, so ist $\Phi_{a,c,\theta} = C$ entweder eine Identität oder eine Transformationsgleichung für $\Psi_{a,c,\theta}$.


\section{Die Transformationsgleichungen für $\gamma_2$ und $\gamma_3$.}\label{sec:71}

Transformationsgleichungen erster Stufe erhält man zunächst aus der Betrachtung der Funktionen \S\ 54, (4), (5):
\[
\gamma_2(\omega) = \sqrt[3]{j(\omega)}, \quad \gamma_3(\omega) = \sqrt{j(\omega) - 1728}.
\]

Wenn $n$ nicht durch 3 teilbar ist, so können wir die Zahlen $a$, $c$, $\theta$ so wählen, daß immer $c$ durch 3 teilbar wird.

Nun übt nach \S\ 54, (15) eine lineare Transformation $\begin{pmatrix} a & b \\ c & d \end{pmatrix}$ auf die Funktion $\gamma_2(\omega)$ den Einfluß:
\[
\gamma_2\left(\frac{a\omega + b}{c\omega + d}\right) = e^{\frac{a+d - (\alpha + \delta)}{3} \cdot 2\pi i} \cdot \gamma_2(\omega).
\]

In den Zusammensetzungen (8), (10) des \S\ 69 wird
\[
\alpha \equiv n, \quad \delta \equiv 0 \pmod{3},
\]
und dann kann man $\beta$ noch so bestimmen, daß auch $a$ und folglich $c_1$ durch 3 teilbar werden.

Daraus ergibt sich, daß durch die beiden Substitutionen
\[
\omega' = \omega + 1, \quad \omega' = -\frac{1}{\omega}
\]
die $\nu$ Funktionen
\begin{equation}\label{eq:71-1}
x_h = e^{\frac{2\pi i h}{3}} \gamma_2\left(\frac{a_h \omega + b_h}{c_h \omega + d_h}\right)
\end{equation}
nur untereinander vertauscht werden und also die Wurzeln einer Transformationsgleichung erster Stufe sind.

Ist zweitens $n$ eine ungerade Zahl, so nehme man $c$ gerade an.

Für die Funktion $\gamma_3(\omega)$ ist [nach \S\ 54, (15)]:
\[
\gamma_3\left(\frac{a\omega + b}{c\omega + d}\right) = e^{\frac{a+d-(\alpha+\delta)}{2} \cdot \pi i} (-1)^{\frac{(c+d)(a+d-1)}{2}} \gamma_3(\omega),
\]
und in den Zusammensetzungen (8), (10), \S\ 69 ist
\[
\alpha \equiv 1, \quad \delta \equiv 0 \pmod{2}, \quad \alpha \equiv 0, \quad \delta \equiv 0 \pmod{2}, \quad \beta \equiv 2\alpha, \quad \gamma \equiv -\alpha \pmod{2}.
\]
Die $\nu$ Größen
\begin{equation}\label{eq:71-2}
x_h = \gamma_3\left(\frac{a_h \omega + b_h}{c_h \omega + d_h}\right)
\end{equation}
vertauschen sich daher untereinander und sind also gleichfalls die Wurzeln einer Transformationsgleichung erster Stufe.

Um für den einfachsten Fall $n = 2$ die erstere dieser Gleichungen zu bilden, beachte man die Relation [\S\ 54, (5)]:
\[
\gamma_2(\omega)^3 - 16 = \frac{f(\omega)^{16} + 16}{f(\omega)^8} = \frac{f(\omega)^{16} + 16}{f(\omega)^8},
\]
woraus wegen
\[
f(\omega_0) f(\omega) = \sqrt[4]{2} \quad [\S\ 34, (16)]
\]
folgt:
\[
\gamma_2(\omega)^3 - 16 = \frac{f(\omega)^{16} + 16}{f(\omega)^8} = f(\omega_0)^8 + \frac{16}{f(\omega)^8}.
\]

Bezeichnen wir diese drei Größen mit $\xi$, $\xi_1$, $\xi_2$, so ergibt sich aus den Relationen (6), (8), \S\ 54
\[
\begin{aligned}
\xi + \xi_1 + \xi_2 &= \gamma_2(\omega)^3, \\
\xi\xi_1 + \xi\xi_2 + \xi_1\xi_2 &= 4 \cdot 9 \cdot \gamma_2(\omega), \\
\xi\xi_1\xi_2 &= -16 \gamma_2(\omega)^2 + 2^6 \cdot 3^3,
\end{aligned}
\]
so daß $\xi$, $\xi_1$, $\xi_2$ die Wurzeln der Gleichung sind
\begin{equation}\label{eq:71-5}
x^3 - \gamma_2(\omega)^2 x^2 + 5 \cdot 9 \cdot \gamma_2(\omega) x + \gamma_2(\omega)^3 - 2^4 \cdot 3^3 \cdot 5^3 = 0.
\end{equation}
Die Gleichung, deren Wurzeln die Kuben der Wurzeln von (5) sind, ist die Invariantengleichung für $n = 2$, und läßt sich daraus ohne Schwierigkeit berechnen.


\section{Multiplikatorgleichungen erster Stufe.}\label{sec:72}

Unter den Multiplikatorgleichungen sollen hier die betrachtet werden, deren Wurzeln die verschiedenen Potenzen von $P_\nu$ sind, multipliziert mit Potenzen von $\gamma_2(\omega)$, $\gamma_3(\omega)$, deren Koeffizienten rational von $j(\omega)$ abhängen. Diese Gleichungen nennen wir \emph{Multiplikatorgleichungen erster Stufe}{}\footnote{Diese Gleichungen sind besonders eingehend von Kiepert untersucht worden, zuerst in mehreren Abhandlungen in Crelles Journal, Bd.~57, 88, 95, am ausführlichsten in den Abhandlungen in den Mathematischen Annalen, Bd.~26, 33. Vgl.\ auch F.~Klein, Mathematische Annalen, Bd.~14, 15 und die oben erwähnten autographierten Vorlesungen.},

Wir betrachten die Größen [\S\ 68, (16)]:
\begin{equation}\label{eq:72-1}
x_{a,c,\theta} = P_{a,c,\theta}^{\frac{c + \theta}{n}} \cdot \gamma_2(\omega)^{-\frac{c}{3}} \cdot \gamma_3(\omega)^{-\frac{c}{2}}
\end{equation}
und den Einfluß, den die Transformationen
\[
\omega' = \omega + 1, \quad \omega' = -\frac{1}{\omega}
\]
auf sie ausüben. Dieser Einfluß bestimmt sich nach den Formeln (8) bis (13), \S\ 69, wonach [weil $c_1 \equiv c \pmod{\theta}$]
\begin{equation}\label{eq:72-2}
\begin{aligned}
x_{a,c,\theta} &\text{ durch } (\omega, \omega+1) \text{ in } e^{\frac{(c+\theta)(A-1)}{12n} \cdot 2\pi i} \cdot x_{a,c,\theta}, \\
&\text{ durch } (\omega, -\tfrac{1}{\omega}) \text{ in } e^{\frac{d-d_1}{2} \cdot \pi i} \cdot x_{a_1,c_1,\theta_1}
\end{aligned}
\end{equation}
übergeht, worin $\varepsilon$ die in \S\ 38, (15) angegebene Bedeutung hat. Es ist aber [\S\ 69, (11), (12)]
\[
d - d_1 = \frac{\alpha + \beta}{\theta} = \frac{1}{\theta}(a - d),
\]
und mithin, wenn wir
\[
\varepsilon\left(\frac{a,c,\theta}{\omega}\right) = e^{\frac{d-d_1}{2} \cdot \pi i} = e^{\frac{a-d}{2\theta} \cdot \pi i} = i^{\frac{a-d}{\theta}}
\]
setzen,
\[
\varepsilon^2 = e^{\frac{a-d}{\theta} \cdot \pi i} = i^{\frac{a-d}{\theta}}
\]
eine 24ste Einheitswurzel, und durch die Vertauschung
\[
\begin{pmatrix} a & b \\ c & d \end{pmatrix} \to \begin{pmatrix} -d & c \\ b & -a \end{pmatrix}
\]
geht
\[
\varepsilon\left(\frac{a,c,\theta}{\omega}\right) \text{ in } \varepsilon\left(\frac{-d,c,\theta}{\omega}\right) = i^{\frac{-d-a}{\theta}} = i^{-\frac{a+d}{\theta}}
\]
über.

Daraus ergibt sich wie oben der Schluß:

\textbf{1.} Für jedes beliebige $n$ sind die Größen
\[
x_{a,c,\theta} = P_{a,c,\theta}^{\frac{24}{n}} \cdot \gamma_2(\omega)^{-\frac{c}{3}} \cdot \gamma_3(\omega)^{-\frac{c}{2}}
\]
die Wurzeln einer Transformationsgleichung.

Ist $n$ durch 3 nicht teilbar, so kann man die $c$, $\theta$, $c_1$ durch 3 teilbar voraussetzen. Dann wird
\[
\alpha \equiv n \equiv \theta \equiv 0, \quad \delta \equiv 0 \pmod{3}.
\]
Es ist also, wie aus \S\ 38, (15) hervorgeht, $\varepsilon$ eine achte Einheitswurzel, beachtet man daher noch
\[
\gamma_2(\omega+1) = \gamma_2(\omega), \quad \gamma_2(-\tfrac{1}{\omega}) = \gamma_2(\omega), \quad \gamma_3(\omega+1) = -\gamma_3(\omega), \quad \gamma_3(-\tfrac{1}{\omega}) = -\gamma_3(\omega),
\]
so haben wir den Satz:

\textbf{2.} Ist $n$ nicht durch 3 teilbar und $c$ durch 3 teilbar, so sind die Größen
\[
P_{a,c,\theta}^{\frac{8}{n}} \cdot \gamma_2(\omega)^{-\frac{c}{3}}
\]
Wurzeln einer Transformationsgleichung.

Um die Anwendung auf den einfachsten Fall $n = 2$ zu machen, setzen wir
\[
x_0 = \frac{f(\omega_0)^4}{f(\omega)^4}, \quad x_1 = \frac{f_1(\omega_0)^4}{f_1(\omega)^4}, \quad x_2 = \frac{f_2(\omega_0)^4}{f_2(\omega)^4},
\]
und dann sind $x_0$, $x_1$, $x_2$ nichts anderes als unsere Funktionen
\[
-f(\omega)^4, \quad f_1(\omega)^4, \quad -f_2(\omega)^4. \quad [\S\ 34, (9)]
\]
Diese sind, wie wir schon früher gesehen haben [\S\ 54, (8)], die Wurzeln der kubischen Gleichung
\begin{equation}\label{eq:72-4}
x^3 - 2\gamma_2(\omega) + 16 = 0.
\end{equation}

Der Umstand, daß $f^4$, $f_1^4$, $-f_2^4$ selbst Wurzeln einer Transformationsgleichung für den Transformationsgrad 2 sind, erklärt die Erscheinung, daß bei Adjunktion dieser Größen, also auch bei Adjunktion von $\varkappa^2$, die zu einem geraden $n$ gehörigen Transformationsgleichungen reducibel werden.

Wenn $n$ eine ungerade Zahl ist, so nehme man $c$ durch 4 teilbar an, dann ist [\S\ 69, (9), (11), (12)]
\[
A \equiv n \pmod{4}, \quad \alpha \equiv 0, \quad \delta \equiv 0, \quad \beta \equiv 2\alpha, \quad \gamma \equiv -\alpha \pmod{4},
\]
also ergibt sich
\[
\varepsilon^2 = i^{\frac{a-d}{\theta}} = i^{\frac{a-d}{\theta}} \cdot (-1)^{\frac{a-d}{2\theta}}
\]
und daraus folgt mit Rücksicht auf
\[
\gamma_3(\omega+1) = -\gamma_3(\omega), \quad \gamma_3(-\tfrac{1}{\omega}) = -\gamma_3(\omega) \quad [\S\ 54, (19)]
\]

\textbf{3.} Ist $n$ ungerade, $c$ durch 4 teilbar, so sind die Größen
\[
P_{a,c,\theta}^{\frac{n}{24}} \cdot \gamma_3(\omega)^{-\frac{c}{2}}
\]
Wurzeln einer Multiplikatorgleichung. Ist $n \equiv 1 \pmod{4}$, so gilt dasselbe von $P_{a,c,\theta}^{\frac{n}{12}}$.

Als Beispiel wählen wir $n = 3$ und setzen
\[
x_0 = \sqrt[3]{\frac{f(\omega_0)}{f(\omega)}} \cdot \gamma_2(\omega)^{-\frac{1}{3}}, \quad x_1 = \sqrt[3]{\frac{f_1(\omega_0)}{f_1(\omega)}} \cdot \gamma_2(\omega)^{-\frac{1}{3}}, \quad x_2 = \sqrt[3]{\frac{f_2(\omega_0)}{f_2(\omega)}} \cdot \gamma_2(\omega)^{-\frac{1}{3}}.
\]

Die Koeffizienten der Gleichung, die diese Größen zu Wurzeln hat, sind rationale Funktionen von $\gamma_2(\omega)$, und da keine der Größen (5) für einen endlichen Wert von $\gamma_2$ unendlich wird, so sind es ganze Funktionen von $\gamma_2$. Nach 3.\ können wir diese Gleichung also in der Form ansetzen:
\[
x^4 + A_1 x^3 + A_2 x^2 + A_3 x + A_4 = 0,
\]
worin $A_1$, $A_2$, $A_3$, $A_4$ ganze rationale Funktionen von $j(\omega)$ sind. Zunächst erhält man $A_4$ als das Produkt $x_0 x_1 x_2 x_3$, welches für keinen Wert von $j(\omega)$ verschwinden kann und daher konstant sein muß. Aus den Anfängen der Entwickelung [\S\ 24]
\[
\begin{aligned}
x_0 &= q^{-\frac{1}{12}} \cdot \sqrt[3]{2} \cdot (1 + q + q^2 + \cdots), \\
x_1 &= q^{\frac{1}{12}} \cdot \sqrt[3]{2} \cdot (1 + q^{\frac{1}{3}} + q^{\frac{2}{3}} + \cdots), \\
x_2 &= q^{\frac{1}{12}} \cdot \sqrt[3]{2} \cdot (1 - q^{\frac{1}{3}} + q^{\frac{2}{3}} - \cdots),
\end{aligned}
\]
findet man daher
\[
A_4 = -27.
\]

Es ist ferner
\[
-A_3 = x_0 + x_1 + x_2 + x_3.
\]
Da die rechte Seite für ein unendliches $\gamma_2$, d.~h.\ für $q = 0$, nicht einmal in der ersten Ordnung unendlich wird, so muß $A_3$ verschwinden.

Aus
\[
-A_2 = x_0 x_1 + x_0 x_2 + x_0 x_3 + x_1 x_2 + x_1 x_3 + x_2 x_3
\]
ergibt sich nach (6) für $A_2$ der konstante Wert 1. Um aber $A_1$ zu bestimmen, müssen wir in der Entwickelung noch um ein Glied weiter gehen. Wir setzen in
\[
A_1 = -(x_0 + x_1 + x_2 + x_3)
\]
für $x_0$ die Entwickelung
\[
x_0 = q^{-\frac{1}{12}} \cdot \sqrt[3]{2} \cdot (1 + q + 5q^2 + \cdots)
\]
und finden $A_1 = 18$, so daß also die gesuchte Gleichung lautet:
\begin{equation}\label{eq:72-7}
x^4 + 18x^2 + 9x - 27 = 0.
\end{equation}

Ist $n$ ungerade und nicht durch 3 teilbar, so nehme man $c$ durch 12 teilbar an.

Es ist alsdann
\[
\alpha \equiv n \equiv \theta \equiv 0, \quad \beta \equiv 2\alpha, \quad \gamma \equiv -\alpha \pmod{12},
\]
und es wird
\[
P_{a,c,\theta}^{\frac{n}{12}} \cdot \gamma_2(\omega)^{-\frac{c}{3}} \cdot \gamma_3(\omega)^{-\frac{c}{2}} = P_{a,c,\theta}^{\frac{n}{12}},
\]
also der Satz

\textbf{4.} Ist $n$ relativ prim zu 6, $c$ durch 12 teilbar, so sind die Größen
\[
P_{a,c,\theta}^{\frac{n}{12}} \cdot \gamma_2(\omega)^{-\frac{c}{3}} \cdot \gamma_3(\omega)^{-\frac{c}{2}}
\]
Wurzeln einer Multiplikatorgleichung erster Stufe.

Da die Funktionen $P^{\frac{n}{12}}$ für keinen endlichen Wert von $\gamma(\omega)$ unendlich oder Null werden, so schließen wir, daß die Koeffizienten der Gleichung, deren Wurzeln die $P^{\frac{n}{12}}$ sind, ganze rationale Funktionen von $\gamma_2$, $\gamma_3$ sind und daß der letzte Koeffizient eine Konstante ist.

Ist $n$ eine Primzahl $p$, so sind die Wurzeln dieser Gleichung
\[
P_{a,c,\theta}^{\frac{p}{12}} = \sqrt[12]{\frac{f(\omega_0)}{f(\omega)}} \cdot \gamma_2(\omega)^{-\frac{c}{3}} \cdot \gamma_3(\omega)^{-\frac{c}{2}}
\]
und die Anfangsglieder der Entwickelungen sind folgende:
\[
x_h = q^{-\frac{p-1}{12p}} \cdot e^{\frac{2\pi i h}{p}} \cdot (1 + \cdots),
\]
wodurch sich für den letzten Koeffizienten der Wert $(-1)^{\frac{p-1}{2}} p$ ergibt. Daraus, daß das erste Glied in der Entwickelung von $j(\omega)$ nach Potenzen von $q$ den Koeffizienten 1 hat, schließt man leicht, daß die numerischen Koeffizienten in diesen Gleichungen rationale ganze Zahlen sind.

Für $p = 5$ hat die fragliche Gleichung die Form:
\[
x^6 + A_1 x^5 + A_2 x^4 + A_3 x^3 + A_4 x^2 + A_5 x + A_6 = 0,
\]
worin die $A_1$, $A_2$, $A_3$, $A_4$ ganze rationale Funktionen von $j(\omega)$ und, wie leicht zu sehen, Konstante sind.

Aus den Anfangsgliedern ergibt sich sofort
\[
A_6 = -5, \quad A_1 = 0, \quad A_5 = -5.
\]
Um aber $A_2$ zu bestimmen, geht man in der Entwickelung von $x_0$ bis zum nächstfolgenden Gliede:
\[
x_0 = q^{-\frac{1}{15}} \cdot (1 + q^{\frac{2}{5}} + q^{\frac{4}{5}} - 2q^{\frac{6}{5}} + \cdots),
\]
woraus man $A_2 = 10$ erhält; also lautet für $n = 5$ die Multiplikatorgleichung
\begin{equation}\label{eq:72-8}
x^6 + 10x^4 - 9x^2 + 5 = 0.
\end{equation}

In gleicher Weise berechnet man die Gleichungen für $p = 7$, $p = 11$, indem man die Entwickelungen von $x_0$ benutzt:
\[
\begin{aligned}
p = 7:\quad &x_0 = q^{-\frac{1}{7}} \cdot (1 + q^{\frac{2}{7}} + q^{\frac{4}{7}} + q^{\frac{6}{7}} + \cdots), \\
p = 11:\quad &x_0 = q^{-\frac{5}{11}} \cdot (1 + q^{\frac{2}{11}} + q^{\frac{4}{11}} + q^{\frac{6}{11}} + q^{\frac{8}{11}} + 2q^{\frac{10}{11}} + \cdots),
\end{aligned}
\]
während von $\gamma_2(\omega)$, $\gamma_3(\omega)$ immer nur die ersten Glieder gebraucht werden. So findet sich
\begin{equation}\label{eq:72-9}
\begin{aligned}
p = 7:\quad &x^8 - 7\gamma_2 x^6 + 7 \cdot 9 \gamma_2^2 x^4 - 7 \cdot 13 \gamma_2^3 x^2 + \gamma_2^4 - 7 = 0, \\
p = 11:\quad &x^{12} - 11\gamma_2 x^{10} + 11 \cdot 40 \gamma_2^2 x^8 + 11 \cdot 153 \gamma_2^3 x^6 \\
&\quad + 11 \cdot 29 \gamma_2^4 x^4 + 99 \gamma_2^5 x^2 - 11 = 0.
\end{aligned}
\end{equation}

Wenn $n$ eine ungerade Quadratzahl ist, so nehmen wir $c$ durch 8 teilbar an. Es sind dann $\alpha:\theta$ und $\delta:\theta$ ebenfalls Quadratzahlen. Es ist ferner
\begin{equation}\label{eq:72-11}
\alpha \equiv \delta \equiv 0 \pmod{8}, \quad A \equiv 1 \pmod{8},
\end{equation}
\[
\alpha:\theta \equiv \alpha, \quad c_1 \equiv \frac{\alpha}{\theta} \pmod{8}, \quad \beta_1 \equiv \frac{\beta}{\theta} \pmod{8}.
\]

Die Einheitswurzel $\varepsilon$ in (2) erhält den Wert
\begin{equation}\label{eq:72-13}
\varepsilon = i^{\frac{1-\beta}{\theta}} = e^{\frac{2\pi i \cdot \alpha}{8}}.
\end{equation}

Es ist aber nach (12) $\beta$ sowohl durch $\theta$ als auch durch $\alpha$ teilbar und $\beta:\theta$ ein Quadrat; also
\[
\varepsilon = i^{\frac{1-\beta}{\theta}} = i^{\frac{1-\beta:\theta}{1}} = e^{\frac{2\pi i \cdot (1-\beta:\theta)}{8}}.
\]
ferner ist nach (12)
\[
\frac{\alpha - \beta}{\theta} = \frac{\alpha}{\theta} - \frac{\beta}{\theta} \equiv \frac{\alpha}{\theta}(1 - \beta:\theta) \pmod{8},
\]
also
\[
\varepsilon = i^{\frac{(\alpha:\theta)(1-\beta:\theta)}{1}} = e^{\frac{2\pi i \cdot (\alpha:\theta)(1-\beta:\theta)}{8}}.
\]

Durch die Vertauschungen (2) geht also
\[
P_{a,c,\theta}^{\frac{n}{24}} \cdot \gamma_2(\omega)^{-\frac{c}{3}} \cdot \gamma_3(\omega)^{-\frac{c}{2}} \text{ in } \varepsilon \cdot P_{a_1,c_1,\theta_1}^{\frac{n}{24}} \cdot \gamma_2(\omega)^{-\frac{c_1}{3}} \cdot \gamma_3(\omega)^{-\frac{c_1}{2}}
\]
über.

Ist $n$ noch durch 3 unteilbar, so nehme man $c$ durch 3 teilbar an, wodurch
\[
\alpha \equiv 1, \quad \delta \equiv 0 \pmod{3}
\]
werden und die in (14) vorkommenden dritten Einheitswurzeln den Wert 1 erhalten.

Hieraus ergeben sich die Sätze:

\textbf{5.} Ist $n$ eine ungerade Quadratzahl, $c$ durch 8 teilbar, so sind die Größen
\[
P_{a,c,\theta}^{\frac{n}{24}} \cdot \gamma_2(\omega)^{-\frac{c}{3}}
\]
und ist $n$ eine durch 3 nicht teilbare ungerade Quadratzahl, $c$ durch 24 teilbar, so sind die Größen
\[
P_{a,c,\theta}^{\frac{n}{24}}
\]
Wurzeln von Multiplikatorgleichungen erster Stufe.

Die ersten Beispiele sind $n = 9$, $n = 25$.

Zur Bildung dieser Gleichungen kann man auf verschiedene Arten gelangen. Wir wollen hier [nach Kiepert{}\footnote{Zur Transformationstheorie der elliptischen Funktionen. Crelles Journal, Bd.~87, 88, 95.}] den Weg gehen, daß wir nach \S\ 69 die Wurzeln einer zum Grade $p$ gehörigen Transformationsgleichung durch die für den Grad $p^2$ rational ausdrücken und diesen Ausdruck in die zu $p$ gehörige Transformationsgleichung einsetzen.

Nach 3., Formel (7) wird die Gleichung
\begin{equation}\label{eq:72-15}
x^4 + 18x^2 + 9\gamma_2(\omega) x - 27 = 0
\end{equation}
von den beiden Funktionen
\[
x = \sqrt[3]{\frac{f(\omega_0)}{f(\omega)}} \cdot \gamma_2(\omega)^{-\frac{1}{3}}
\]
befriedigt. Bezeichnen wir den ersten dieser Werte mit $x$, so geht der zweite durch die Substitution $\begin{pmatrix} 0 & -1 \\ 1 & 0 \end{pmatrix}$ über, und folglich wird durch $x$ auch die Gleichung
\begin{equation}\label{eq:72-17}
x^4 + 18 \cdot 27 x^2 - \gamma_2(\omega) x - 27 = 0
\end{equation}
befriedigt. Setzen wir nun
\begin{equation}\label{eq:72-18}
y = \sqrt[3]{\frac{f(\frac{\omega}{3})}{f(3\omega)}} \cdot \gamma_2(3\omega)^{-\frac{1}{3}},
\end{equation}
so geht $x$ durch die Substitution $\begin{pmatrix} 0 & -1 \\ 3 & 0 \end{pmatrix}$ in $-\frac{1}{y}$ über, so daß man auch die Gleichung erhält:
\begin{equation}\label{eq:72-19}
y^4 + 18y^2 - \gamma_2(3\omega) y - 27 = 0,
\end{equation}
und durch Elimination von $\gamma_2(3\omega)$ aus (17) und (19) erhält man nach Weghebung des Faktors $y^2 - 27$
\begin{equation}\label{eq:72-20}
x^4 - 18x^2y^2 - y^2(y^4 - 27y^2 + 27^2) = 0.
\end{equation}
Löst man diese quadratische Gleichung nach $x^2$ auf, so folgt:
\begin{equation}\label{eq:72-21}
x^2 = y^4 + 9y^2 + 27 = (y^2 + 3)^3 - 27,
\end{equation}
wo über das Zeichen durch Einsetzen der Anfangsglieder der Entwickelungen entschieden wird, am einfachsten wohl, da diese Gleichung (nach \S\ 69) auch für
\[
y = \sqrt[3]{j(3\omega)}, \quad x = \sqrt[3]{j(\omega)}
\]
erfüllt wird, indem man
\[
x = y^2 + 3
\]
setzt.

Sondert man in (15) die erste Potenz von $x$ ab und erhebt ins Quadrat, so erhält man durch Einsetzen von (21) die gesuchte Multiplikatorgleichung für den 9ten Transformationsgrad. Sie erhält eine einfachere Gestalt, wenn man
\begin{equation}\label{eq:72-22}
t = x^2 + 9 = y^4 + 9y^2 + 36
\end{equation}
setzt:
\begin{equation}\label{eq:72-23}
[\gamma_2(\omega)^2 - 27 \cdot 64](t - 27) = (t^2 - 36t + 27 \cdot 8)^2
\end{equation}
oder
\begin{equation}\label{eq:72-24}
\gamma_2(\omega)^2 (t - 27) = t^3 - 24t^2.
\end{equation}

Ganz ähnlich kann man beim 25sten Transformationsgrad verfahren.

Wenn wir
\[
x = \sqrt[6]{\frac{f(\omega_0)}{f(\omega)}} \cdot \gamma_2(\omega)^{-\frac{1}{3}}
\]
setzen, so haben wir zunächst nach 4.\ (8):
\begin{equation}\label{eq:72-26}
x^6 + 10x^4 - \gamma_2(\omega) x^2 + 5 = 0,
\end{equation}
woraus, wie oben, die beiden Gleichungen
\[
\begin{aligned}
x^6 + 10 \cdot 5^2 x^4 - \gamma_2(5\omega) x^2 + 5^5 &= 0, \\
y^6 + 10y^4 - \gamma_2(5\omega) y^2 + 5 &= 0,
\end{aligned}
\]
und durch Elimination von $\gamma_2(5\omega)$
\[
x^6 - 10x^2y^2(5 + y^2) - y^2(y^6 + 5y^4 + 52y^2 + 51y^2 + 5^2).
\]
Diese Gleichung nach $x^2$ aufgelöst, ergibt:
\begin{equation}\label{eq:72-27}
x^2 = y^6 + 5y^4 + 15y^2 + 25 = (y^2 + 1)^5 + 25,
\end{equation}
und wenn wir zur Abkürzung die rechte Seite dieser Gleichung mit $\chi(y)$ bezeichnen, nach (26)
\begin{equation}\label{eq:72-28}
\chi(y)^3 + 10\chi(y)^2 - \gamma_2(\omega) \chi(y) + 5^2 = 0.
\end{equation}
Dies ist die gesuchte Gleichung 30sten Grades für $y$.


\section{Die Schlaeflischen Modulargleichungen.}\label{sec:73}

Zu einfacheren Transformationsgleichungen gelangt man, wenn man dem Rationalitätsbereich, der bis jetzt aus den rationalen Funktionen von $j(\omega)$ bestand, die Größe $f(\omega)^2$ adjungiert. Diesem Rationalitätsbereich gehören die Funktionen von $\omega$ an, die durch die beiden Substitutionen
\begin{equation}\label{eq:73-1}
\omega' = \omega + 2, \quad \omega' = -\frac{1}{\omega}
\end{equation}
ungeändert bleiben (\S\ 54, 2). Wenn also ein System von $\nu$ Funktionen $\Phi_{a,c,\theta}$ durch die Substitutionen (1) nur unter sich permutiert wird, so sind diese Funktionen die Wurzeln einer Transformationsgleichung, deren Koeffizienten rational von $f(\omega)^2$ abhängen. Hierzu gehören (wie wir früher schon auf anderem Wege nachgewiesen haben) für ein ungerades $n$, das wir jetzt immer voraussetzen, gewisse Potenzen der Größen
\[
f\left(\frac{c + \omega}{\theta}\right)^2
\]
worin, wie ein für allemal bemerkt sei, $c$ durch 16, und wenn $n$ nicht durch 3 teilbar ist, durch 48 teilbar angenommen wird.

Die etwas erweiterten Grundsätze des \S\ 70 führen verhältnismäßig einfach zur Berechnung dieser Gleichungen.

Eine Erweiterung ist aber notwendig aus folgendem Grunde:

Bei den bisherigen Betrachtungen konnten wir den Schluß machen: wenn eine rationale Funktion von $j(\omega)$ für jedes endliche $\omega$ mit positiv imaginärem Teil endlich bleibt, so ist sie eine ganze Funktion von $j(\omega)$, weil zu jedem endlichen $j(\omega)$ auch ein endliches, nicht reelles $\omega$ gehört (\S\ 52). Bei den rationalen Funktionen von $f(\omega)$ können wir aber aus der Endlichkeit für jedes endliche imaginäre $\omega$ nur schließen, daß sie ganze Funktionen von $f(\omega)$ und $1:f(\omega)$ sind, weil nur zu jedem endlichen $f(\omega)$ mit Ausnahme von $f(\omega) = 0$ ein endliches imaginäres $\omega$ gehört.

Es entspricht aber der Substitution
\begin{equation}\label{eq:73-2}
\omega' = -\frac{1}{\omega}
\end{equation}
die Vertauschung
\[
f(\omega) \to f_2(\omega), \quad f_1(\omega) \to f(\omega), \quad f_2(\omega) \to f_1(\omega) \quad [\S\ 34, (13)]
\]
und wenn also die Funktionen $\Phi_{a,c,\theta}$ so gewählt sind, daß sie auch durch die Substitution (2) nur untereinander permutiert werden, so werden die Koeffizienten der Gleichung, deren Wurzeln sie sind, rational von
\[
f(\omega)^2 + \frac{16}{f(\omega)^2}
\]
abhängen. Sie sind ganze Funktionen dieser Verbindungen, wenn die $\Phi_{a_h,c_h,\theta_h}$ für jedes endliche, nicht reelle $\omega$ endlich bleiben, und sie sind konstant, wenn alle $\Phi_{a_h,c_h,\theta_h}$ auch für $q = 0$, d.~h.\ für $f(\omega) = 0$ und $f(\omega) = \infty$, endlich bleiben. In diesem Falle sind sämtliche $\Phi_{a_h,c_h,\theta_h}$ einer und derselben Konstanten gleich (\S\ 70).

Daraus ergibt sich der Satz:

\begin{satz}
Bildet man ganze rationale Funktionen $\Phi_{a_h,c_h,\theta_h}$ aus
\[
f\left(\frac{c + \omega}{\theta}\right)^2, \quad f_1\left(\frac{c + \omega}{\theta}\right)^2, \quad \frac{1}{f\left(\frac{c + \omega}{\theta}\right)^2}, \quad \frac{1}{f_1\left(\frac{c + \omega}{\theta}\right)^2},
\]
welche die Eigenschaft haben:
\begin{enumerate}[label=\arabic*.]
\item durch die Substitutionen
\[
\omega' = \omega + 2, \quad \omega' = \omega - 2, \quad \omega' = -\frac{1}{\omega}
\]
nur untereinander vertauscht zu werden,
\item für $q = 0$ nicht unendlich zu werden,
\end{enumerate}
so ist
\[
\Phi_{a,c,\theta} = \text{const.}
\]
eine Transformationsgleichung.
\end{satz}

Um diese Bedingungen zu befriedigen, ist zunächst der Einfluß der Substitutionen (1) auf die Funktionen $f$ zu untersuchen. Dieser ergibt sich aus den Zusammensetzungen \S\ 69, (8) bis (13) und aus den Transformationsformeln für die $f$-Funktionen \S\ 40.

Zur Abkürzung führen wir die Bezeichnung ein:
\begin{equation}\label{eq:73-3}
\begin{aligned}
u &= f_0(\omega)^2, &\quad v &= f_0\left(\frac{c + \omega}{\theta}\right)^2, \\
\nu_1 &= f_1(\omega)^2, &\quad v_1 &= f_1\left(\frac{c + \omega}{\theta}\right)^2, \\
\nu_2 &= f_2(\omega)^2, &\quad v_2 &= f_2\left(\frac{c + \omega}{\theta}\right)^2.
\end{aligned}
\end{equation}

Es ergeben sich dann folgende zusammengehörige Änderungen, wenn in der Bezeichnung auf die Verwandlung der Zahlen $a$, $c$, $\theta$ in $a_1$, $c_1$, $\theta_1$ oder $A_1$, $C_1$, $\theta_1$, wie sie eben durch die angeführten Formeln des \S\ 69 charakterisiert ist, keine Rücksicht genommen wird:

\begin{center}
\begin{tabular}{c|ccc}
& $\nu$ & $v$ & $v_1$ \\
\midrule
$\omega \to \omega + 2$ & $\nu$ & $v$ & $v_1$ \\
$\omega \to \omega - 2$ & $\nu$ & $v$ & $v_1$ \\
$\omega \to -\frac{1}{\omega}$ & $\nu_2$ & $v_2$ & $v$ \\
\end{tabular}
\end{center}

worin $\varrho_1$, $\varrho_2$ zwei dritte Einheitswurzeln sind, die, falls $n$ nicht durch 3 teilbar ist, den Wert 1 haben.

Um ferner die Wirkung der Substitution (2), die aus der Transformation zweiten Grades
\[
\omega' = -\frac{1}{\omega}
\]
hervorgeht, auf die Funktionen $u$, $v$ zu ermitteln, müssen wir die Transformationen erster und $n$ter Ordnung
\[
\begin{pmatrix} a & b \\ c & d \end{pmatrix}, \quad \begin{pmatrix} \alpha & \beta \\ 0 & \theta \end{pmatrix}
\]
so bestimmen, daß
\[
\begin{pmatrix} a & b \\ c & d \end{pmatrix} \cdot \begin{pmatrix} 1 & 0 \\ 0 & n \end{pmatrix} = \begin{pmatrix} \alpha & \beta \\ 0 & \theta \end{pmatrix} \cdot \begin{pmatrix} a_1 & b_1 \\ c_1 & d_1 \end{pmatrix}
\]
wird. Dieser Ansatz führt zu den Gleichungen
\[
\theta = \alpha\delta - \beta\gamma + \theta(\alpha'\delta' - \beta'\gamma'), \quad \alpha = \delta\theta + \beta\theta', \quad \nu\alpha = \theta a_1, \quad \nu\beta = \theta b_1.
\]
Hiermit ist zunächst, da $\alpha + \beta$ und $\gamma - \delta$ zufolge $\alpha\delta - \beta\gamma = 1$ relativ prim sind, $\theta'$ bestimmt als der größte gemeinschaftliche Teiler von $\alpha$ und $c + \delta$, und aus $n = \alpha'\theta'$ ergibt sich $\alpha'$. Dann ist nach den Gleichungen (6):
\[
\alpha + \beta = \frac{c}{\theta'}, \quad \gamma = \frac{c'}{\theta'}, \quad \delta = \frac{c''}{\theta'},
\]
und $\delta$ und $\beta$ lassen sich so bestimmen, daß
\begin{equation}\label{eq:73-7}
(\alpha + \beta)\delta - \beta(\gamma + \delta) = 1.
\end{equation}

Aus den Gleichungen (6) für $\alpha$ und $c - \delta$ findet sich dann
\begin{equation}\label{eq:73-8}
c' + \alpha = \alpha\delta - (\alpha + \beta)(\gamma - \delta).
\end{equation}
$\delta$ und $\beta$ können, da $\alpha - \beta$ und $\gamma + \delta$ beide ungerade sind, nach (7) nicht beide ungerade sein; folglich fällt $c'$ nach (8) gerade aus. Ersetzt man, was nach (7) gestattet ist, $\delta$, $\beta$ durch $\delta + h(\gamma - \delta)$, $\beta + h(\alpha + \beta)$, für ein beliebiges $h$, so ändert sich $c'$ nach (6) und (8) um $2h\alpha'$, und über $h$ kann so verfügt werden, daß $c'$ durch 16 oder 48 teilbar wird.

Es ist dann nach (6) $\alpha + \beta \equiv \gamma - \delta$ gerade und daher die Formel (12), \S\ 40 anzuwenden. Darin ist nun zu berücksichtigen
\[
\alpha - \beta \equiv \alpha, \quad \gamma - \delta \equiv c \pmod{16}, \quad \alpha + \beta \equiv \nu, \quad \delta(\gamma - \delta) \equiv -\theta \pmod{16},
\]
woraus folgt:
\[
n(\alpha - \beta)^2 \equiv \nu^2 + \theta \equiv \nu - \theta' \pmod{16},
\]
also
\[
\left(\frac{\alpha - \beta}{2}\right)^2 \equiv \frac{\nu - \theta'}{n} \pmod{4}
\]
und
\[
\left(\frac{\gamma - \delta}{2}\right)^2 \equiv \frac{(\alpha + \beta) + (\gamma - \delta) - \theta'}{n} \cdot \theta' \equiv \frac{c' - \theta'}{n} \pmod{4}.
\]

Ist $n$ nicht durch 3 teilbar, so ist
\[
\alpha + \beta = \alpha'(1 + 0), \quad \gamma - \delta = \alpha'(1 - 0), \quad 2\beta = \alpha'(1 - \alpha'), \quad 2\gamma = 2\alpha'(1 - \alpha'),
\]
also entweder $\alpha$ und $\delta$ oder $\beta$ und $\gamma$ durch 3 teilbar und also
\[
\varrho_1 = \varrho_2 = 1.
\]

Hieraus ergeben sich folgende zusammengehörige Vertauschungen (wobei die Änderung von $a$, $c$, $\theta$ in $a'$, $c'$, $\theta'$ durch (6) bestimmt ist):

\begin{center}
\begin{tabular}{c|ccc}
& $u$ & $\nu$ & $v$ \\
\midrule
$\omega \to \omega + 2$ & $u$ & $\nu$ & $v$ \\
$\omega \to -\frac{1}{\omega}$ & $u'$ & $\nu'$ & $v'$ \\
\end{tabular}
\end{center}

Wir wenden die Vertauschungen (4), (5), (9) auf folgende Funktionen an:
\begin{equation}\label{eq:73-10}
\begin{aligned}
A &= u^r v^s + u^s v^r, \\
B &= u^r u_1^s v^s + u^s u_1^r v^r,
\end{aligned}
\end{equation}
worin $r$, $s$ zwei ganze Zahlen sind, die den Bedingungen
\begin{equation}\label{eq:73-11}
(n-1)r \equiv 0, \quad (n+1)s \equiv 0 \pmod{12}
\end{equation}
genügen, die also, wenn $n$ durch 3 teilbar ist, beide durch 3 teilbar sind und
\[
\frac{r \cdot s}{\varrho_1 \cdot \varrho_2}
\]
ist.

Es ergeben sich dann folgende zusammengehörige Vertauschungen:

\begin{center}
\begin{tabular}{c|cc}
& $A$ & $B$ \\
\midrule
$\omega \to \omega + 2$ & $A$ & $B$ \\
$\omega \to -\frac{1}{\omega}$ & $\varrho_1^{\frac{(n-1)r}{12}} A$ & $\varrho_1^{\frac{(n-1)r}{12}} \varrho_2^{\frac{(n+1)s}{12}} B$ \\
\end{tabular}
\end{center}

Bilden wir nun aus $A$, $B$ eine ganze rationale Funktion mit numerischen Koeffizienten $M_{h,k}$
\begin{equation}\label{eq:73-15}
\Phi_{a,c,\theta} = \sum_{h,k} M_{h,k} A^h B^k,
\end{equation}
worin, falls in (12) Vorzeichenänderungen auftreten, die Exponenten $h$, $k$ so einzurichten sind, daß in allen Gliedern von (13) die gleichen Vorzeichenveränderungen stattfinden, so wird das Funktionensystem $\Phi_{a_h,c_h,\theta_h}$ oder wenigstens $\Phi_{a_h,c_h,\theta_h}^3$ der Forderung 1.\ des oben aufgestellten Satzes genügen und wir haben, um auch die Forderung 2.\ zu befriedigen und so eine Transformationsgleichung zu erhalten, die Koeffizienten $M_{h,k}$ so zu bestimmen, daß die sämtlichen $\nu$ Werte von $\Phi_{a_h,c_h,\theta_h}$ für $q = 0$ endlich bleiben. Diese Aufgabe vereinfacht sich wesentlich, wenn $n$ keinen quadratischen Teiler hat, und noch mehr, wenn $n$ eine Primzahl ist.

Hat nämlich $n$ keinen quadratischen Teiler, so kann man aus $\Phi_{1,0,n}$ die sämtlichen Werte $\Phi_{a_h,c_h,\theta_h}$ herleiten, durch Vermehrung von $\omega$ um gewisse ganze Zahlen; wenn also, was unsere Forderung ist, in der Entwickelung von $\Phi_{1,0,n}$ nach steigenden Potenzen von $q$ keine negativen Potenzen vorkommen, so gilt das Gleiche von sämtlichen $\Phi_{a_h,c_h,\theta_h}$.

Ist aber $n$ eine Primzahl, so genügt der Nachweis, daß $\Phi_{1,0,n}$ keine negativen Potenzen von $q$ enthält, da das Gleiche durch Vertauschung von $n$ mit $\omega:n$ für $\Phi_{\omega,0,1}$ folgt.

Der konstante Wert, den die Funktion $\Phi_{a_h,c_h,\theta_h}$ erhält, bestimmt sich aus einem Gliede der Entwickelung.

Bei der Ausführung dieser Rechnungen dienen die Entwickelungen \S\ 24, (11),
\[
\begin{aligned}
u &= q^{-\frac{1}{12}} \cdot 2^{\frac{1}{3}} \cdot (1 + q + 2q^2 + \cdots)^{\frac{2}{3}}, \\
\nu_1 &= 2q^{\frac{1}{12}} \cdot (1 + q^{\frac{2}{3}} + q^{\frac{4}{3}} + \cdots)^{\frac{2}{3}}, \\
v &= q^{-\frac{n}{12(n+1)}} \cdot 2^{\frac{1}{3}} \cdot (1 + q^{\frac{2n}{n+1}} + \cdots)^{\frac{2}{3}}, \\
v_1 &= 2q^{\frac{n}{12(n+1)}} \cdot (1 + q^{\frac{2}{n+1}} + \cdots)^{\frac{2}{3}}.
\end{aligned}
\]

\begin{equation}\label{eq:73-14}
\begin{aligned}
u &= q^{-\frac{1}{12}} \cdot 2^{\frac{1}{3}} \cdot (1 - q^{2r} + q^{4r} - \cdots)^{\frac{1}{3}} (1 + q^{2r} + q^{4r} + \cdots)^{\frac{1}{3}}, \\
\nu_1 &= q^{\frac{1}{12}} \cdot 2^{\frac{1}{3}} \cdot (1 + q^{2r} + q^{4r} + \cdots)^{\frac{1}{3}} (1 + q^{2r} + q^{4r} + \cdots)^{\frac{1}{3}}, \\
v &= q^{-\frac{n}{12(n+1)}} \cdot 2^{\frac{1}{3}} \cdot (1 - q^{\frac{2n}{n+1}} + q^{\frac{4n}{n+1}} - \cdots)^{\frac{1}{3}} (1 + q^{\frac{2n}{n+1}} + \cdots)^{\frac{1}{3}}, \\
v_1 &= q^{\frac{n}{12(n+1)}} \cdot 2^{\frac{1}{3}} \cdot (1 + q^{\frac{2}{n+1}} + q^{\frac{4}{n+1}} + \cdots)^{\frac{1}{3}} (1 + q^{\frac{2}{n+1}} + \cdots)^{\frac{1}{3}}.
\end{aligned}
\end{equation}

Die Formeln werden nicht immer am einfachsten, wenn $r$, $s$ möglichst klein angenommen werden, sondern es erweist sich am zweckmäßigsten, noch die Bedingung
\[
\frac{(n-1)r}{12} = \frac{(n+1)s}{12}
\]
hinzuzufügen.

Wir erhalten so für die sieben ersten ungeraden Primzahlen folgende Bestimmung von $A$ und $B$:

\begin{center}
\begin{tabular}{c|cc}
$n$ & $A$ & $B$ \\
\midrule
$3$ & $\left(\frac{u}{v}\right)^2 + \left(\frac{v}{u}\right)^2$ & $\left(\frac{uv_1}{u_1v}\right)^2 + \left(\frac{u_1v}{uv_1}\right)^2$ \\
$5$ & $\left(\frac{u}{v}\right)^2 + \left(\frac{v}{u}\right)^2$ & $\left(\frac{uv_1}{u_1v}\right)^2$ \\
$7$ & $\left(\frac{u}{v}\right)^2$ & $\left(\frac{uv_1}{u_1v}\right)^2$ \\
$11$ & $\left(\frac{u}{v}\right)^2$ & $\left(\frac{u_1}{v_1}\right)^2$ \\
$13$ & $\left(\frac{u}{v}\right)^2$ & $\left(\frac{u_1}{v_1}\right)^2$ \\
$17$ & $u^2 v^2$ & $u_1^2 v_1^2$ \\
$19$ & $u^2 v^2$ & $u_1^2 v_1^2$ \\
\end{tabular}
\end{center}

und daraus erhält man durch Elimination der negativen Potenzen die gesuchten Gleichungen zwischen $A$ und $B$:

\begin{center}
\begin{tabular}{c|l}
$n$ & Gleichung \\
\midrule
$3$ & $A^3 - B^2 - 54 = 0$ \\
$5$ & $A^2 - B - 22 = 0$ \\
$7$ & $A^3 - B^2 + 2A - B = -1$ \\
$11$ & $A^5 - B^3 + 17AB - 54A^2 + 53B + 116A - 44 = 0$ \\
$13$ & $A^4 - B^3 + 19AB^2 - 95A^2B - 109A^3$ \\
& $\quad - 25B^2 + \cdots = 0$ \\
\end{tabular}
\end{center}

\footnotetext{Diese Gleichungen sind zuerst von Schlaefli aufgestellt (Journal für Mathematik, Bd.~72).}

Aus diesem System von Gleichungen leitet man ein zweites und drittes her für die Funktionen $f_1$, $f_2$, indem man $\omega$ durch $\omega + 1$ und darauf $\omega$ durch $-\frac{1}{\omega}$ ersetzt. Diese beiden Systeme haben die gleiche Form, nur ist das eine Mal
\[
u = f(\omega)^2, \quad \nu_1 = f_1(\omega)^2, \quad \nu_2 = f_2(\omega)^2
\]
das andere Mal
\[
u = f_1(\omega)^2, \quad \nu_1 = f_2(\omega)^2, \quad \nu_2 = f(\omega)^2
\]
zu setzen. Aus den Vertauschungen (4) ergibt sich so:

\begin{center}
\begin{tabular}{c|ll}
$n=3$ & $A_1 = \left(\frac{u_1}{v_1}\right)^2 + \left(\frac{v_1}{u_1}\right)^2$ & $B_1 = \left(\frac{u_1v_2}{u_2v_1}\right)^2 + \left(\frac{u_2v_1}{u_1v_2}\right)^2$ \\
& $A_1^3 - B_1^2 - 54 = 0$ & \\
\midrule
$n=5$ & $A_1 = \left(\frac{u_1}{v_1}\right)^2 + \left(\frac{v_1}{u_1}\right)^2$ & $B_1 = \left(\frac{u_1v_2}{u_2v_1}\right)^2$ \\
& $A_1^2 - B_1 - 22 = 0$ & \\
\midrule
$n=7$ & $A_1 = \left(\frac{u_1}{v_1}\right)^2$ & $B_1 = \left(\frac{u_1v_2}{u_2v_1}\right)^2$ \\
& $A_1^3 - A_1 - B_1 = 2$ & \\
\end{tabular}
\end{center}

\begin{center}
\begin{tabular}{c|ll}
$n=11$ & $A_1 = \left(\frac{u_1}{v_1}\right)^2$ & $B_1 = \left(\frac{u_2}{v_2}\right)^2$ \\
& $A_1^5 - B_1^3 + 17A_1B_1 - 34A_1^2 - 34B_1 + 116A_1$ & \\
& $\quad - 44 = 0$ & \\
\midrule
$n=19$ & $A_1 = \left(\frac{u_1}{v_1}\right)^2$ & $B_1 = \left(\frac{u_2}{v_2}\right)^2$ \\
& $A_1^5 - B_1^2 - 9A_1^3B_1 - 9A_1B_1^3 + 26A_1^4$ & \\
& $\quad + 26B_1^4 - 128B_1 - 128A_1 = 0$ & \\
\end{tabular}
\end{center}


\section{Die Form der Schlaeflischen Modulargleichungen für einen Primzahlgrad.}\label{sec:74}

Die Form, die wir im vorigen Paragraphen für die zwischen
\begin{equation}\label{eq:74-1}
u = f_0(\omega), \quad v = f_0\left(\frac{c + \omega}{\theta}\right)
\end{equation}
bestehenden Relationen gefunden haben, läßt sich, wenigstens wenn der Transformationsgrad eine Primzahl $p$ ist, leicht unter ein allgemeines Gesetz bringen. Da für die Folge viel auf diese Form ankommt, gehen wir hier noch etwas genauer darauf ein.

Die Bestimmung der Zahlen $r$, $s$ nach (11) und (15) des vorigen Paragraphen hängt von dem Verhalten von $p$ gegen den Modul 24 ab, und da wir den Fall $p = 5$ ausschließen können, so haben folgende Fälle:
\begin{equation}\label{eq:74-2}
\begin{array}{c|c|c}
p \pmod{24} & r & s \\
\hline
p \equiv 1 & \frac{p-1}{12} & \frac{p-1}{12} \\
p \equiv 5 & \frac{p+1}{12} & \frac{p+1}{12} \\
p \equiv 7 & \frac{p+1}{12} & \frac{p-1}{12} \\
p \equiv 11 & \frac{p+1}{12} & \frac{p+1}{12} \\
p \equiv 13 & \frac{p-1}{12} & \frac{p-1}{12} \\
p \equiv 17 & \frac{p-1}{12} & \frac{p+1}{12} \\
p \equiv 19 & \frac{p+1}{12} & \frac{p+1}{12} \\
p \equiv 23 & \frac{p+1}{12} & \frac{p+1}{12} \\
\end{array}
\end{equation}
so daß $r - s$ stets ungerade ist und $p-1$ durch $2r$, $p-1$ durch $2s$ teilbar ist. Hiernach wird
\begin{equation}\label{eq:74-3}
\begin{aligned}
A &= \left(\frac{u}{v}\right)^r + \left(\frac{v}{u}\right)^r, \\
B &= \left(\frac{uv_1}{u_1v}\right)^s + \left(\frac{u_1v}{uv_1}\right)^s = \frac{(uv_1)^s + (u_1v)^s}{(uv)^s}.
\end{aligned}
\end{equation}

Nun wissen wir, daß die $p-1$ Größen $v$ die Wurzeln einer irreducibeln Transformationsgleichung
\begin{equation}\label{eq:74-4}
\Phi_p(v, u) = v^{p+1} - U_1 v^p + \cdots + U_{p+1} = 0
\end{equation}
sind, in der die Koeffizienten $U_1$, \ldots, $U_{p+1}$ rationale Funktionen von $u$ sind.

Wir schließen sofort, daß es ganze rationale Funktionen von $u$ sind. Denn erstens werden für $u = 0$ die sämtlichen Wurzeln von (4) unendlich, wie man erkennt, wenn man $\omega = i\infty$, also $q \to 0$ werden läßt. Zweitens geht nach (9) des vorigen Paragraphen durch die Vertauschung $\begin{pmatrix} 0 & -1 \\ 1 & 0 \end{pmatrix}$ die Gesamtheit der Wurzeln $v$ in
\[
v' = \frac{1}{v} \cdot \frac{u}{u'}
\]
über. Hieraus folgt, daß für $u = 0$ die sämtlichen Wurzeln $v$ in Null übergehen, also keine von ihnen unendlich wird, woraus zu schließen ist, daß nicht nur die $U_1$, $U_2$, \ldots\ $U_{p+1}$ ganze Funktionen von $u$ sind, sondern daß auch jede von ihnen den Faktor $u$ enthalten muß.

Wir schließen nun zunächst, genau wie bei der Invariantengleichung (\S\ 69, 2., 3.\ und 6.), daß
\begin{equation}\label{eq:74-5}
\Phi_p(v, u) = \Phi_p(u, v),
\end{equation}
und daß
\begin{equation}\label{eq:74-6}
\Phi_p(v, u) = (v^{p+1} - u^{p+1}) - p \sum_{h,k=1}^{p} a_{h,k} v^h u^k,
\end{equation}
worin $a_{h,k} = a_{p+1-h,\,p+1-k}$ ganze Zahlen sind.

Wir wollen diese Funktion in der Weise darstellen
\begin{equation}\label{eq:74-7}
\Phi_p(v, u) = v^{p+1} - u^{p+1} - p \sum_{h,k=1}^{p} a_{h,k} v^h u^k,
\end{equation}
worin also $a_{h,k} = a_{p+1-h,\,p+1-k}$ ebenfalls ganze Zahlen sind, auf deren Teilbarkeit durch $p$ es nun weiter nicht ankommt, und es ist insbesondere $a_{1,0} = -1$. Wenn wir in der Gleichung
\[
\Phi_p[f(p\omega), f(\omega)] = 0
\]
$\omega$ durch $\omega + 2$ ersetzen, so ergibt sich wegen der Irreducibilität auf Grund der Relation
\[
f(\omega + 2)^2 = e^{\frac{2\pi i}{12}} f(\omega)^2 = \varrho \cdot f(\omega)^2;
\]
daß in (7) nicht alle Glieder, sondern nur solche vorkommen, in denen die Exponenten $h$, $k$ der Kongruenz
\begin{equation}\label{eq:74-8}
hp + k \equiv p + 1 \pmod{24}
\end{equation}
genügen.

Nun kennen wir noch eine weitere Eigenschaft der Funktionen $\Phi_p(v, u)$, die sich aus der schon benutzten Vertauschung (9) des vorigen Paragraphen ergibt, wo $\theta = 1$ zu setzen ist, und
\[
v \to \frac{1}{v} \cdot u
\]
die, wenn wir zur Abkürzung
\[
\varepsilon = e^{\frac{2\pi i}{24}}
\]
setzen, so dargestellt werden kann:
\begin{equation}\label{eq:74-9}
\Phi_p(v, u) = \left(\frac{v}{u}\right)^{p+1} \Phi_p\left(\frac{u}{v} \cdot u, v\right).
\end{equation}

Hieraus schließt man auf die Relationen
\begin{equation}\label{eq:74-10}
\varepsilon^{h+k-p-1} \cdot a_{p+1-h,\,p+1-k} = a_{h,k} = \varepsilon^{h+k-p-1} \cdot a_{h,k}.
\end{equation}

Wenn wir also in (7) die Glieder mit gleichen Koeffizienten $a_{h,k}$ zusammenfassen, so können wir uns auf die Annahme beschränken, daß $h > k$, $h + k \geq p + 1$ sei, und es ergibt sich $\Phi_p$, von den beiden Gliedern $v^{p+1}$, $u^{p+1}$ abgesehen, als ein Aggregat von Gliedern von den folgenden beiden Formen (wenn noch berücksichtigt wird, daß wegen (8) $h \equiv k$):
\begin{equation}\label{eq:74-11}
v^h u^k + v^k u^h = v^{\frac{h+k-p-1}{2}} (v^{k+1} u^{\frac{h-k}{2}} + v^{\frac{h-k}{2}} u^{k+1}) \cdot (vu)^{\frac{h+k-p-1}{2}},
\end{equation}
\begin{equation}\label{eq:74-12}
v^h u^k - v^k u^h = \varepsilon^{\frac{h-k}{2}} \cdot v^{\frac{h+k-p-1}{2}} (v^{k+1} u^{\frac{h-k}{2}} - v^{\frac{h-k}{2}} u^{k+1}) \cdot (vu)^{\frac{h+k-p-1}{2}}.
\end{equation}

Die Koeffizienten dieser Glieder in $\Phi_p$ sind ganze Zahlen.

Nun ergibt sich aus (8):
\[
h + k - p - 1 \equiv -h(p-1) \pmod{24},
\]
und daher ist $h - k$ durch $2r$, $h + k - p - 1$ (worin $h$ auch $= k$ sein kann) durch $2s$ teilbar [\S\ 74, (2)].

Setzen wir daher
\[
\frac{h - k}{2r} = \alpha, \quad \frac{h + k - p - 1}{2s} = \beta,
\]
so sind $\alpha$ und $\beta$ positive ganze Zahlen, die zwischen den Grenzen
\begin{equation}\label{eq:74-13}
0 \leq \alpha \leq \frac{p-1}{2r}, \quad 0 \leq \beta \leq \frac{p+1}{2s}
\end{equation}
liegen, und überdies gilt die aus (8) folgende Kongruenz $-h(p-1) = 2s\beta \pmod{24}$, daß wenigstens in den Fällen, wo $\varepsilon = -1$ ist, $k \equiv \beta \pmod{2}$ [\S\ 74, (2)], also stets $h \equiv k$.

Demnach wird (11)
\[
(v^r)^\alpha (u^r)^\alpha \cdot [(vu)^s]^\beta + [(vu)^s]^\beta (v^r)^\alpha (u^r)^\alpha
\]
und (12)
\[
(v^r)^\alpha (u^r)^\alpha \cdot [(vu)^s]^\beta - [(vu)^s]^\beta (v^r)^\alpha (u^r)^\alpha.
\]

Nun gelten die bekannten Formeln, wenn $x$, $y$ beliebige Größen sind und
\[
z = x + y
\]
gesetzt wird,
\[
x^n + y^n = z^n - n z^{n-2} xy + \frac{n(n-3)}{2} z^{n-4} x^2 y^2 - \cdots,
\]
\[
(x^n - y^n)(x - y) = z^{n+1} - \frac{n+1}{1} z^{n-1} xy + \frac{(n+1)n}{1 \cdot 2} z^{n-3} x^2 y^2 - \cdots,
\]
woraus man durch den Schluß von $n$ auf $n-1$ erkennt, daß
\[
\frac{x^n - y^n}{x - y}
\]
sich für jedes beliebige $n$ als ganze rationale Funktion $n$ten Grades von $z$ darstellen läßt, die, wenn $z$ eine ganze Zahl ist, ganzzahlige Koeffizienten hat, deren höchster $=1$ ist.

Die beiden Größen
\[
\left(\frac{v}{u}\right)^r + \left(\frac{u}{v}\right)^r \quad \text{und} \quad (vu)^s - \frac{1}{(vu)^s}
\]
können also in dieser Weise als ganze rationale Funktionen von $A$ und $B$ der Grade $\alpha$ und $\beta$ dargestellt werden, und wir finden, wenn wir noch das dem Werte $\beta = \frac{p+1}{2s}$ entsprechende Glied, das den Koeffizienten $-1$ hat, absondern und mit $c_{\alpha,\beta}$ ganzzahlige Koeffizienten bezeichnen:

\begin{equation}\label{eq:74-14}
\Phi_p(v, u) = v^{p+1} + u^{p+1} + \sum_{\alpha,\beta} c_{\alpha,\beta} A^\alpha B^\beta,
\end{equation}
worin $\alpha$ und $\beta$ an die Grenzen
\[
0 \leq \alpha \leq \frac{p-1}{2r}, \quad 0 \leq \beta \leq \frac{p+1}{2s}
\]
gebunden sind. Dies ist die Form, die wir im vorigen Paragraphen den Modulargleichungen bis $p = 19$ gegeben haben.


\section{Die irrationalen Formen der Modulargleichungen.}\label{sec:75}

Es ist noch zu erwähnen, daß die in \S\ 69, 4., 5.\ für die Invariantengleichung bei zusammengesetztem Transformationsgrade durchgeführte Betrachtung unverändert auch für die Schlaeflischen Modulargleichungen gilt, woraus wir schließen können, daß alle diese Gleichungen rationale Zahlenkoeffizienten haben.

Den Transformationsgleichungen lassen sich durch Anwendung desselben Verfahrens weit einfachere Formen geben. Die Gleichungsformen, mit denen wir uns jetzt beschäftigen werden, enthalten die drei Funktionen $f$, $f_1$, $f_2$ zugleich; da man aber nach \S\ 34, (11), (12) zwei dieser Funktionen durch die dritte ausdrücken kann, so lassen sich zwei von ihnen eliminieren, und man kann so zu den Gleichungen des vorigen Paragraphen gelangen. Wenn man für $f_1$, $f_2$ die Ausdrücke durch $f$, oder für die drei Funktionen $f$, $f_1$, $f_2$ die Ausdrücke durch $\varkappa^2$ einsetzt, so kommen Wurzelzeichen vor, woraus sich der Name dieser Gleichungen erklärt.

Wir setzen wie oben
\begin{equation}\label{eq:75-1}
\begin{aligned}
u &= f(\omega)^2, &\quad v &= f\left(\frac{c + \omega}{\theta}\right)^2, \\
\nu_1 &= f_1(\omega)^2, &\quad v_1 &= f_1\left(\frac{c + \omega}{\theta}\right)^2, \\
\nu_2 &= f_2(\omega)^2, &\quad v_2 &= f_2\left(\frac{c + \omega}{\theta}\right)^2,
\end{aligned}
\end{equation}
und erhalten folgende zusammengehörige Vertauschungen:

\begin{center}
\begin{tabular}{c|cccc}
& $\nu$ & $v$ & $\nu_1$ & $v_1$ \\
\midrule
$\omega \to \omega + 1$ & $\varrho \nu$ & $\varrho v$ & $\varrho_1 \nu_1$ & $\varrho_1 v_1$ \\
$\omega \to \omega - 1$ & $\varrho^{-1} \nu$ & $\varrho^{-1} v$ & $\varrho_1^{-1} \nu_1$ & $\varrho_1^{-1} v_1$ \\
$\omega \to -\frac{1}{\omega}$ & $\nu_2$ & $v_2$ & $\nu$ & $v$ \\
\end{tabular}
\end{center}
worin $\varrho$, $\varrho_1$ die oben [\S\ 73, (4)] definierten dritten Einheitswurzeln sind.

Wir unterscheiden drei verschiedene Fälle nach dem Verhalten von $n$ zum Modul 8.

\textbf{1.} $n + 1 \equiv 0 \pmod{8}$.

Wir setzen
\[
\begin{aligned}
A &= uv - \frac{(-1)^{\frac{n+1}{8}}}{uv} \cdot \nu_1 v_1, \\
B &= uv \cdot u_1 v_1 + \frac{(-1)^{\frac{n+1}{8}}}{uv \cdot u_1 v_1} \cdot (uv)^2,
\end{aligned}
\]
\begin{equation}\label{eq:75-3}
\begin{aligned}
A &= uv - \frac{(-1)^{\frac{n+1}{8}}}{uv} \cdot \nu_1 v_1, \\
B &= uv \cdot u_1 v_1 + \frac{(-1)^{\frac{n+1}{8}}}{uv \cdot u_1 v_1} \cdot (uv)^2, \\
C &= \frac{u}{v} + \frac{v}{u} + \frac{u_1}{v_1} + \frac{v_1}{u_1},
\end{aligned}
\end{equation}
so daß sich die zusammengehörigen Vertauschungen ergeben

\begin{center}
\begin{tabular}{c|ccc}
& $A$ & $B$ & $C$ \\
\midrule
$\omega \to \omega + 2$ & $A$ & $B$ & $C$ \\
$\omega \to -\frac{1}{\omega}$ & $A$ & $B$ & $C$ \\
$\omega \to \omega - 1$ & $\varrho^{\frac{(n+1)r}{12}} A$ & $\varrho^{\frac{(n+1)r}{12}} B$ & $\varrho^{\frac{(n+1)r}{12}} C$ \\
\end{tabular}
\end{center}

Ein Produkt von der Form $A^h B^k$ nimmt also durch die beiden Vertauschungen
\[
\omega' = \omega + 2, \quad \omega' = \omega - 1
\]
die Faktoren an
\[
\varrho^{\frac{(n+1)(h-k)r}{12}}, \quad \varrho^{\frac{(n+1)(h-k)r}{12}},
\]
die $= 1$ sind, wenn $n - 1$ durch 3 teilbar ist. Wir bilden also jetzt mit numerischen Koeffizienten $M_{h,k}$ Funktionen der Form
\begin{equation}\label{eq:75-5}
\Phi_{a,c,\theta} = \sum_{h,k} M_{h,k} A^h B^k,
\end{equation}
worin, wenn $n \equiv 0$ oder $\equiv 1 \pmod{3}$ ist, $h$ und $k$ nur solche (ganzzahlige) Werte annehmen dürfen, deren Differenzen $h - k$ bei der Teilung mit 3 denselben Rest lassen, wenn aber $n \equiv -1 \pmod{3}$ dieser Beschränkung nicht unterworfen sind. Eine solche Funktion selbst, oder wenigstens ihre dritte Potenz genügt also einer Transformationsgleichung erster Stufe. Sie bleibt außerdem für alle endlichen Werte von $j(\omega)$ endlich und ist folglich eine ganze algebraische Funktion von $\gamma(\omega)$. [Die Transformation zweiter Ordnung $\begin{pmatrix} 0 & -1 \\ 1 & 0 \end{pmatrix}$ braucht hier nicht zugezogen zu werden, da der Inbegriff der $\Phi_{a_h,c_h,\theta_h}$ schon bei der Substitution $(\omega, \omega + 1)$, nicht erst bei $(\omega, \omega - 2)$ ungeändert bleibt. Vgl.\ die Bemerkung am Anfang des \S\ 73.]

Wenn wir also die Konstanten in $\Phi_{a_h,c_h,\theta_h}$ so bestimmen, daß in den Entwickelungen dieser Funktionen keine negativen Potenzen von $q$ vorkommen, so müssen die sämtlichen $\Phi_{a_h,c_h,\theta_h}$ einer und derselben Konstanten gleich sein.

Zur Erreichung dieses Zieles genügt es auch hier, wenn $n$ keinen quadratischen Teiler hat, daß $\Phi_{1,0,n}$, und wenn $n$ eine Primzahl ist, wenn $\Phi_{1,0,n}$ für $q = 0$ endlich bleibt.

Bei diesen Rechnungen machen wir Gebrauch von den Entwickelungen:
\[
\begin{aligned}
u^r &= q^{-\frac{r}{12}} \cdot 2^{\frac{r}{3}} \cdot \prod_{m=1}^{\infty} (1 + q^{2m-1})^{\frac{2r}{3}} (1 - q^{2m-1})^{\frac{2r}{3}}, \\
\nu_1^s &= 2^s q^{\frac{s}{12}} \cdot \prod_{m=1}^{\infty} (1 + q^{2m})^{\frac{2s}{3}},
\end{aligned}
\]
\begin{equation}\label{eq:75-6}
\begin{aligned}
u^r &= q^{-\frac{r}{12}} \cdot 2^{\frac{r}{3}} (1 + rq^2 + \cdots), \\
v^r &= q^{-\frac{nr}{12(n+1)}} \cdot 2^{\frac{r}{3}} (1 + rq^{\frac{2n}{n+1}} + \cdots), \\
\nu_1^s &= 2^s q^{\frac{s}{12}} (1 + sq^2 + \cdots), \\
v_1^s &= 2^s q^{\frac{ns}{12(n+1)}} (1 + sq^{\frac{2}{n+1}} + \cdots).
\end{aligned}
\end{equation}

Die letzteren Formeln sind richtig für $n > 7$ (für $n = 3, 5, 7$ sind die Glieder von $q^3$, $q^5$, $q^7$ an zu modifizieren). Für $n = 7$, $n = 23$ zeigen diese Ausdrücke, daß $A$ selbst für $q = 0$ endlich bleibt, woraus für diese Fälle die Gleichungen folgen:
\begin{equation}\label{eq:75-8}
\begin{aligned}
n = 7:\quad &A^2 - B = -1, \\
n = 23:\quad &A^3 - B^2 + 2A - B = -1.
\end{aligned}
\end{equation}

Durch einfache Rechnung findet man ferner noch:
\begin{equation}\label{eq:75-9}
\begin{aligned}
n = 3:\quad &A^2 - B - A = 0, \\
n = 11:\quad &A^3 - A - B = 2, \\
n = 19:\quad &A^3 - A^2 + 2A - B = 1.
\end{aligned}
\end{equation}

Nicht ganz so einfach gestaltet sich die Rechnung für ein zusammengesetztes $n$. So muß man z.~B.\ für $n = 15$ die Bedingung der Endlichkeit für $q = 0$ nicht nur für $\Phi_{1,0,15}$, sondern auch für $\Phi_{3,0,5}$ berücksichtigen. Diese beiden aber genügen. Man erhält so:
\begin{equation}\label{eq:75-10}
n = 15:\quad A^3 - AB + 11 = 0.
\end{equation}

\textbf{2.} Ist $n \equiv 3 \pmod{8}$, so sind dieselben Schlüsse zu ziehen, wenn wir setzen:
\begin{equation}\label{eq:75-11}
\begin{aligned}
A &= uv - u_1v_1 - uv_1, \\
B &= uv \cdot u_1v_1 + uv \cdot u_1v_1 - uv \cdot u_1v_1,
\end{aligned}
\end{equation}
für die man aus (2) die zusammengehörigen Vertauschungen erhält:

\begin{center}
\begin{tabular}{c|cc}
& $A$ & $B$ \\
\midrule
$\omega \to \omega + 2$ & $A$ & $B$ \\
$\omega \to -\frac{1}{\omega}$ & $\varrho^{\frac{(n+1)r}{12}} A$ & $\varrho^{\frac{(n+1)r}{12}} B$ \\
$\omega \to \omega - 1$ & $\varrho^{\frac{(n+1)r}{12}} A$ & $\varrho^{\frac{(n+1)r}{12}} B$ \\
\end{tabular}
\end{center}

und hieraus leitet man in der gleichen Weise die Gleichungen ab:
\begin{equation}\label{eq:75-13}
\begin{aligned}
n = 3:\quad &A^2 - B - A = 0, \\
n = 11:\quad &A^5 - A^3 - B^2 + A = 0, \\
n = 19:\quad &A^3 - 7A^2 - B = 1.
\end{aligned}
\end{equation}

\textbf{3.} Ist $n \equiv 1 \pmod{4}$, so kann man ebenso verfahren mit den Funktionen
\begin{equation}\label{eq:75-14}
\begin{aligned}
A &= uv + u_1v_1, \\
B &= uv \cdot u_1v_1 + uv \cdot u_1v_1 - uv \cdot u_1v_1,
\end{aligned}
\end{equation}
für die man die zusammengehörigen Vertauschungen erhält:

\begin{center}
\begin{tabular}{c|cc}
& $A$ & $B$ \\
\midrule
$\omega \to \omega + 2$ & $A$ & $B$ \\
$\omega \to -\frac{1}{\omega}$ & $A$ & $B$ \\
$\omega \to \omega - 1$ & $\varrho^{\frac{(n+1)r}{12}} A$ & $\varrho^{\frac{(n+1)r}{12}} B$ \\
\end{tabular}
\end{center}

Nur der erste Fall $n = 5$ führt hier zu einem einfachen Resultat:
\begin{equation}\label{eq:75-16}
n = 5:\quad A^2 - B = 0.
\end{equation}

Für $n = 5$ läßt sich noch eine einfachere Form der Transformationsgleichung gewinnen.

Wenn wir nämlich auf die drei Funktionen
\[
w = \frac{u^2 + uv}{v^2}, \quad w_1 = \frac{u_1^2 + u_1v_1}{v_1^2}, \quad w_2 = \frac{u_2^2 + u_2v_2}{v_2^2}
\]
die Substitutionen $(\omega, \omega)$; $(\omega, \omega + 1)$ anwenden, so ergeben sich unter der Voraussetzung $n \equiv 5 \pmod{8}$ die zusammengehörigen Vertauschungen:

\begin{center}
\begin{tabular}{c|ccc}
& $w$ & $w_1$ & $w_2$ \\
\midrule
$\omega \to \omega + 2$ & $w$ & $w_1$ & $w_2$ \\
$\omega \to \omega - 1$ & $\varrho^{\frac{n-5}{12}} w$ & $\varrho^{\frac{7(n-5)}{12}} w_1$ & $\varrho^{\frac{7(n-5)}{12}} w_2$ \\
\end{tabular}
\end{center}

Setzt man also
\[
\begin{aligned}
A &= w^2 + w_1^2 + w_2^2, \\
B &= ww_1 + ww_2 + w_1w_2, \\
C &= ww_1w_2,
\end{aligned}
\]
so erhält man

\begin{center}
\begin{tabular}{c|ccc}
& $A$ & $B$ & $C$ \\
\midrule
$\omega \to \omega + 2$ & $A$ & $B$ & $C$ \\
$\omega \to \omega - 1$ & $A$ & $B$ & $C$ \\
\end{tabular}
\end{center}

so daß zwei dieser Funktionen ebenso wie oben $A$ und $B$ benutzt werden können. Für $n = 5$ zeigt sich aber, daß in den Entwickelungen von $w$, $w_1$, $w_2$ nach Potenzen von $q$ keine negativen Potenzen vorkommen und daß also $A$, $B$, $C$ und mithin auch $w$, $w_1$, $w_2$ selbst konstant sind.

Es läßt sich also die Modulargleichung für $n = 5$ in jeder der drei Formen aufstellen:
\begin{equation}\label{eq:75-18}
\begin{aligned}
u^2 v^2 - v^2 u^2 &= 2uv, \\
u v^2 - v u^2 &= 2uv, \\
u v^2 - v^2 u &= 2uv.
\end{aligned}
\end{equation}
Diese Gleichungen lassen sich auch aus der von Jacobi (Fund.\ art.\ 30, gesammelte Werke, Bd.~1, S.~123) gegebenen herleiten.

Wir schließen diese Betrachtungen, indem wir in den einfachsten Fällen die Jacobischen Modulargleichungen aus diesen irrationalen Formen ableiten.

Wir setzen [vgl.\ \S\ 54, (3)]:
\[
\begin{aligned}
x &= \sqrt[4]{\frac{u_1}{v_1}} = \sqrt[4]{\frac{f_1(\omega)^2}{f_1(\frac{c+\omega}{\theta})^2}}, &\quad y &= \sqrt[4]{\frac{u}{v}} = \sqrt[4]{\frac{f(\omega)^2}{f(\frac{c+\omega}{\theta})^2}}, \\
x' &= \sqrt[4]{\frac{v_1}{u_1}} = \sqrt[4]{\frac{f_1(\frac{c+\omega}{\theta})^2}{f_1(\omega)^2}}, &\quad y' &= \sqrt[4]{\frac{v}{u}} = \sqrt[4]{\frac{f(\frac{c+\omega}{\theta})^2}{f(\omega)^2}},
\end{aligned}
\]
und eliminieren mittels der Relationen
\begin{equation}\label{eq:75-19}
\begin{aligned}
x^4 + \frac{1}{\varkappa^2} \cdot \varkappa'^2 &= 1, &\quad y^4 + \frac{1}{\varkappa^2} \cdot \varkappa'^2 &= 1, \\
\frac{1}{\varkappa^2} &= \frac{u^4}{16}, &\quad \frac{1}{\varkappa'^2} &= \frac{v^4}{16},
\end{aligned}
\end{equation}
die Größen $u$ und $v^2$.

Für $n = 3$ erhält man aus (13):
\begin{equation}\label{eq:75-20}
xy^3 = -1,
\end{equation}
eine Form der Modulargleichung, die von Legendre herrührt.

Daraus, indem man für $x^4$, $y^4$ aus (19) die Werte setzt,
\[
\varkappa^2 \varkappa'^2 + y^4 - 4x^2 y^2 - \varkappa^2 \varkappa'^2 x^4 y^4 = 4\varkappa^2 \varkappa'^2 x^2 y^2 (1 - x^2 y^2)^2
\]
oder
\[
(x^2 - y^2)^2 = 4(1 - x^2 y^2)^2 \varkappa^2 \varkappa'^2 x^2 y^2
\]
und indem man hieraus die Wurzel zieht und das Vorzeichen durch $q = 0$ bestimmt, findet man:
\begin{equation}\label{eq:75-21}
x^2 - y^2 = 2xy \sqrt{1 - x^2 y^2} \cdot \varkappa \varkappa',
\end{equation}
was, abgesehen von dem dort nicht näher erklärten Vorzeichen von $\sqrt{A}$, mit \S\ 10, (12) übereinstimmt.

Für $n = 5$ folgt aus der zweiten Gleichung (18):
\[
(x^2 - y^2)^2 = 4xy(1 - x^2 y^2),
\]
und daraus durch Quadrieren
\[
(x^2 - y^2)^4 = 16x^2 y^2 (1 - x^2 y^2)^2,
\]
was leicht in die Form gebracht wird
\[
(x^2 - y^2)^2 = 4xy(1 - x^2 y^2)(1 + x^2 y^2).
\]
Zieht man hieraus die Wurzel, so folgt, wie oben:
\begin{equation}\label{eq:75-22}
x^2 - y^2 = 2xy \sqrt{(1 - x^2 y^2)(1 + x^2 y^2)} \quad (n = 5).
\end{equation}

Für $n = 7$ erhält man, ohne Wurzelziehen, aus (8):
\begin{equation}\label{eq:75-23}
x^4 y^4 - xy = 1,
\end{equation}
und daraus:
\begin{equation}\label{eq:75-24}
\begin{aligned}
&x^8 + y^8 - 8xy(1 + x^4 y^4) + 28x^2 y^2 (1 + 2x^4 y^4) \\
&\quad - 56x^3 y^3 (1 + x^4 y^4) + 70x^4 y^4 = 0 \quad (n = 7).
\end{aligned}
\end{equation}


\section{Zusammengesetzte Transformationsgrade.}\label{sec:76}

Ist der Transformationsgrad eine zusammengesetzte Zahl, so kann man noch einfachere Transformationsgleichungen aufstellen als die, die man auf dem Wege des vorigen Paragraphen gewinnt.

Wir führen diese Betrachtungen hier nur in den einfachsten Fällen durch.

Der ungerade Transformationsgrad $n$ sei in zwei Faktoren zerlegt
\begin{equation}\label{eq:76-1}
n = n' \cdot n'',
\end{equation}
die zueinander relativ prim sind.

Es lassen sich dann jeder Transformation $n$ter Grades von der Form
\begin{equation}\label{eq:76-2}
\omega' = \frac{a\omega + b}{c\omega + d}
\end{equation}
je eine und nur eine Transformation der Grade $n'$, $n''$:
\begin{equation}\label{eq:76-3}
\omega' = \frac{a'\omega + b'}{c'\omega + d'}, \quad \omega' = \frac{a''\omega + b''}{c''\omega + d''}
\end{equation}
zuordnen, die durch folgende Bedingungen bestimmt sind:
\begin{equation}\label{eq:76-4}
\theta' c' \equiv c \pmod{\alpha'}, \quad \theta'' c'' \equiv c \pmod{\alpha''}, \quad c'd'' = c \pmod{\alpha''}, \quad c''d' = c \pmod{\alpha'},
\end{equation}
und umgekehrt folgt aus jedem Paar Transformationen von der Form (3) nach (4) eine und nur eine Transformation (2). Nach (4) sind nämlich zunächst $\theta'$, $\theta''$ bestimmt als die größten gemeinschaftlichen Teiler von $c$ mit $n'$ und $n''$, und darauf wird $c'$ nach dem Modul $\alpha'$, $c''$ nach dem Modul $\alpha''$ bestimmt aus den beiden letzten Kongruenzen (4).

Es kommt nun vor allem darauf an, zu zeigen, daß die Kongruenzen (4) erhalten bleiben, wenn die Transformationen (2) und (3) nach \S\ 69, (8) bis (12) durch die beiden linearen Transformationen
\[
\omega' = \omega - 1, \quad \omega' = \omega + 1
\]
umgeformt werden.

Wir setzen nach den erwähnten Formeln:
\[
\begin{pmatrix} a_1 & b_1 \\ c_1 & d_1 \end{pmatrix} = \begin{pmatrix} a & b \\ c & d \end{pmatrix} \cdot \begin{pmatrix} 1 & -1 \\ 0 & 1 \end{pmatrix}, \quad \begin{pmatrix} a_2 & b_2 \\ c_2 & d_2 \end{pmatrix} = \begin{pmatrix} a & b \\ c & d \end{pmatrix} \cdot \begin{pmatrix} 1 & 1 \\ 0 & 1 \end{pmatrix},
\]
\begin{equation}\label{eq:76-5}
\begin{pmatrix} a_1 & b_1 \\ c_1 & d_1 \end{pmatrix} = \begin{pmatrix} a' & b' \\ c' & d' \end{pmatrix} \cdot \begin{pmatrix} 1 & -1 \\ 0 & 1 \end{pmatrix}, \quad \begin{pmatrix} a_2 & b_2 \\ c_2 & d_2 \end{pmatrix} = \begin{pmatrix} a'' & b'' \\ c'' & d'' \end{pmatrix} \cdot \begin{pmatrix} 1 & 1 \\ 0 & 1 \end{pmatrix}.
\end{equation}

Es folgt zunächst aus \S\ 69, (9):
\[
c' = c + \theta' d \pmod{\alpha'}, \quad c'' = c' d'' + \theta'' \pmod{\alpha''}, \quad c'd'' = c'' + \theta'' \pmod{\alpha''},
\]
woraus, nach (4)
\[
\theta' c'_1 = c'_1 d'_1 + \theta'_1 = c_1 \pmod{\alpha'}, \quad \theta'' c''_1 = c''_1 d''_1 + \theta''_1 = c_1 \pmod{\alpha''}
\]
in Übereinstimmung mit den Kongruenzen (4).

Für die Zusammensetzung (6) ergibt sich nach \S\ 69, (12), daß $\theta$, $\theta'$, $\theta''$ die größten gemeinschaftlichen Teiler von $a$, $c$; $a'$, $c'$; $a''$, $c''$ sind; weil aber $\theta''$ relativ prim zu $a'$, und $\theta'' c'' \equiv c \pmod{\alpha'}$ ist, so ist auch $\theta$ der größte gemeinschaftliche Teiler von $a'$ und $c$ und aus den gleichen Gründen $\theta'$ der größte gemeinschaftliche Teiler von $a''$ und $c$, woraus man schließt, da $a'$, $a''$ relativ prim sind:
\[
\theta = \theta' \cdot \theta''.
\]

Ferner ist nach \S\ 69, (11), (12), (13):
\[
\beta = \alpha\delta - \beta\gamma = -\alpha\theta, \quad \gamma = \nu\alpha' - \beta\delta' = -\beta\theta \cdot \frac{\alpha'}{\theta'},
\]
also nach (4):
\[
\beta\alpha' - \delta\alpha = \alpha\alpha'(d'd'' - cc') = 0 \pmod{\alpha'},
\]
folglich auch
\[
\beta\theta - \gamma\theta = 0 \pmod{\alpha'},
\]
in Übereinstimmung mit (4), und ebenso folgt:
\[
\gamma\theta' - \delta\theta = 0 \pmod{\alpha''},
\]
wodurch also der Beweis geführt ist, daß die durch (4) ausgedrückte Zusammengehörigkeit der Transformationen
\[
\begin{pmatrix} a & b \\ c & d \end{pmatrix}, \quad \begin{pmatrix} a' & b' \\ c' & d' \end{pmatrix}, \quad \begin{pmatrix} a'' & b'' \\ c'' & d'' \end{pmatrix}
\]
durch Anwendung irgend einer linearen Transformation auf $\omega$ nicht gestört wird. Da wir hier $n$ als ungerade voraussetzen, so können wir immer $c$, $c'$, $c''$ durch 16 teilbar annehmen; und wenn $n$ und folglich auch $n'$, $n''$ durch 3 unteilbar sind, so können $c$, $c'$, $c''$ auch durch 3 teilbar angenommen werden. Ist aber $n$ durch 3 teilbar, so wird von den beiden Faktoren $n'$, $n''$ der eine, etwa $n''$, durch 3 teilbar sein, der andere, $n'$, nicht. Es kann dann $c'$ noch durch 3 teilbar vorausgesetzt werden, nicht aber $c$ und $c''$. In diesem Falle soll die Abhängigkeit des $c''$ von $c$ noch näher bestimmt werden durch die Kongruenz
\begin{equation}\label{eq:76-7}
\theta'' c'' \equiv c \pmod{3\alpha''}.
\end{equation}

Eine Lösung dieser Kongruenz kann man immer aus einer Lösung der Kongruenz (4) $\theta'' c'' \equiv c \pmod{\alpha''}$ herleiten, indem man zu $c''$ ein Vielfaches von $\alpha''$ hinzufügt.

Die Kongruenz (7) hat dann nach \S\ 69, (9) bis (13) zur Folge:
\begin{equation}\label{eq:76-8}
\beta \equiv n\alpha' \pmod{3}, \quad \alpha \equiv \alpha'' \cdot \theta' \cdot \theta'' \pmod{3},
\end{equation}
\begin{equation}\label{eq:76-9}
\nu\alpha' \equiv \alpha\delta' - \beta\gamma' \pmod{3}.
\end{equation}

Es mögen nun $v$, $v_1$, $v_2$; $v'$, $v'_1$, $v'_2$; $v''$, $v''_1$, $v''_2$ dieselbe Bedeutung für die Zahlen $n'$, $n''$ haben, welche den $v$, $v_1$, $v_2$ in \S\ 73, (3) für die Zahl $n$ gegeben war, nämlich:
\[
\begin{aligned}
v &= f(\omega)^2, &\quad v_1 &= f_1(\omega)^2, &\quad v_2 &= f_2(\omega)^2, \\
v' &= f(n'\omega)^2, &\quad v'_1 &= f_1(n'\omega)^2, &\quad v'_2 &= f_2(n'\omega)^2, \\
v'' &= f(n''\omega)^2, &\quad v''_1 &= f_1(n''\omega)^2, &\quad v''_2 &= f_2(n''\omega)^2.
\end{aligned}
\]

Wir wenden die Vertauschungstabelle (4), \S\ 73 auf diese Funktionen an. Da $n'$ unter allen Umständen durch 3 unteilbar ist, so sind die kubischen Einheitswurzeln $\varrho'$, $\varrho'' = 1$ zu setzen, während infolge der Kongruenzen (8), (9):
\begin{equation}\label{eq:76-11}
\varrho'' = \varrho^{n'}, \quad \varrho = \varrho''^{n'}
\end{equation}
wird. Wir erhalten hiernach folgende zusammengehörige Vertauschungen:

\begin{center}
\begin{tabular}{c|cccccc}
& $v$ & $v_1$ & $v_2$ & $v'$ & $v'_1$ & $v'_2$ \\
\midrule
$\omega \to \omega + 2$ & $v$ & $v_1$ & $v_2$ & $v'$ & $v'_1$ & $v'_2$ \\
$\omega \to \omega - 1$ & $\varrho v$ & $\varrho^2 v_1$ & $v_2$ & $\varrho v'$ & $\varrho^2 v'_1$ & $v'_2$ \\
$\omega \to -\frac{1}{\omega}$ & $v_2$ & $v$ & $v_1$ & $v''_2$ & $v''$ & $v''_1$ \\
\end{tabular}
\end{center}

\begin{equation}\label{eq:76-12}
\begin{tabular}{c|cccccc}
& $v$ & $v_1$ & $v_2$ & $v''$ & $v''_1$ & $v''_2$ \\
\midrule
$\omega \to \omega + 2$ & $v$ & $v_1$ & $v_2$ & $v''$ & $v''_1$ & $v''_2$ \\
$\omega \to \omega - 1$ & $\varrho v$ & $\varrho^2 v_1$ & $v_2$ & $\varrho^{n'} v''$ & $\varrho^{2n'} v''_1$ & $v''_2$ \\
$\omega \to -\frac{1}{\omega}$ & $v_2$ & $v$ & $v_1$ & $v''_2$ & $v''$ & $v''_1$ \\
\end{tabular}
\end{equation}
die zusammen mit den A [\S\ 73, (4)] gelten.

Wir unterscheiden jetzt zwei Fälle.

\textbf{1.} Wenn
\begin{equation}\label{eq:76-13}
(n'+1)n' + (n''+1)n'' \equiv 0 \pmod{8}
\end{equation}
ist, so setzen wir
\begin{equation}\label{eq:76-14}
\begin{aligned}
U &= uv \cdot u'v', \quad V = uv \cdot u''v'', \quad W = u'v' \cdot u''v'', \\
A &= U + V + W, \\
B &= UV + UW + VW.
\end{aligned}
\end{equation}

Aus (12) und \S\ 73, (4) erhalten wir dann folgende zusammengehörige Vertauschungen:

\begin{center}
\begin{tabular}{c|cc}
& $A$ & $B$ \\
\midrule
$\omega \to \omega + 2$ & $A$ & $B$ \\
$\omega \to -\frac{1}{\omega}$ & $A$ & $B$ \\
$\omega \to \omega - 1$ & $\varrho^{\frac{(n'+1)r'}{12}} A$ & $\varrho^{\frac{(n'+1)r'}{12}} B$ \\
\end{tabular}
\end{center}

Wir wenden unser Prinzip zur Herleitung von Modulargleichungen auf diese Funktionen an und bemerken dazu noch folgendes:

Die in (16) vorkommenden dritten Einheitswurzeln sind $= 1$, wenn entweder $n$ durch 3 unteilbar und $n'$ durch 3 teilbar ist, oder $n''$ durch 3 teilbar ist und $n'$ den Rest 2 läßt. In diesen Fällen ist jede rationale Funktion von $A$ und $B$ Wurzel einer Transformationsgleichung. In den anderen Fällen kommt diese Eigenschaft dem Kubus einer solchen rationalen Funktion von $A$ und $B$ zu, bei denen die Differenzen der Exponenten sämtlicher Glieder einander nach dem Modul 3 kongruent sind.

Sind diese rationalen Funktionen ganze Funktionen und sind außerdem ihre sämtlichen Werte für $q = 0$ endlich, so müssen sie einer Konstanten gleich sein, und dadurch gewinnen wir Transformationsgleichungen.

Sind $n'$, $n''$ Primzahlen, so genügt es auch hier (vgl.\ \S\ 72), wenn die negativen Potenzen von $q$ in der Entwickelung einer solchen Funktion nach steigenden Potenzen von $q$ in dem einen Hauptfall wegfallen, nämlich in dem, wo
\[
\begin{pmatrix} a & b \\ c & d \end{pmatrix}, \quad \begin{pmatrix} a' & b' \\ c' & d' \end{pmatrix}, \quad \begin{pmatrix} a'' & b'' \\ c'' & d'' \end{pmatrix}
\]
gleich sind
\[
\begin{pmatrix} 1 & 0 \\ 0 & n \end{pmatrix}, \quad \begin{pmatrix} 1 & 0 \\ 0 & n' \end{pmatrix}, \quad \begin{pmatrix} 1 & 0 \\ 0 & n'' \end{pmatrix},
\]
also
\[
U = f(\omega)^2 f(n'\omega)^2 f(n''\omega)^2 f(nn'\omega)^2.
\]

\end{document}
